\documentclass[11pt]{amsart}
\usepackage{amssymb,amsmath,booktabs,longtable}
\usepackage[hidelinks,
  pdftitle={New lower bounds for kissing numbers in dimensions 25 through 96},
  pdfauthor={Alexey Kravatskiy}]{hyperref}

\newtheorem{theorem}{Theorem}[section]
\newtheorem{lemma}[theorem]{Lemma}
\newtheorem{proposition}[theorem]{Proposition}
\newtheorem{corollary}[theorem]{Corollary}
\theoremstyle{remark}
\newtheorem{remark}[theorem]{Remark}
\newtheorem{example}[theorem]{Example}

\newcommand{\Leech}{\Lambda_{24}}
\newcommand{\Pff}{P_{48}}
\newcommand{\Gsz}{\Gamma_{72}}
\newcommand{\ip}[2]{\langle #1,#2\rangle}
\newcommand{\R}{\mathbb{R}}
\newcommand{\Z}{\mathbb{Z}}
\newcommand{\Q}{\mathbb{Q}}

\title[Kissing numbers in dimensions 25 to 96]
      {New lower bounds for kissing numbers\\in dimensions 25 through 96}
\author{Alexey Kravatskiy}
\address{MIRIAI (Moscow Independent Research Institute\newline
of Artificial Intelligence), Moscow, Russia}
\email{kravatskii.a@miriai.org}
\date{30 August 2026}
\subjclass[2020]{Primary 52C17; Secondary 11H31, 11H71, 05B40, 94B25}
\keywords{Kissing number, spherical code, extremal even unimodular lattice,
spherical design, Leech lattice, linear programming bound, exact verification}

\raggedbottom

\begin{document}

\begin{abstract}
The kissing number $K(n)$ is the largest number of unit vectors in $\R^n$ with
pairwise inner products at most $\tfrac12$. We improve the best known lower
bound in $47$ dimensions between $25$ and $96$. All but four of these bounds
come from the extremal even unimodular lattices $\Leech$, $\Pff$ and $\Gsz$, by
two constructions: cross-sections of the minimal shell and layered
configurations built on it.

For cross-sections we apply Venkov's theorem, that the minimal shell of an
extremal even unimodular lattice of dimension divisible by $24$ is a spherical
$11$-design, to monomials in the inner products with several fixed minimal
vectors. The resulting linear constraints on the $k$-point distribution,
together with two consequences of the lattice structure, bound the number of
minimal vectors orthogonal to all of them. Each bound carries an exact rational
dual certificate. Which Gram matrix to fix is settled by Hermite's inequality,
and the matrices attaining its bound are exhibited in $\Gsz$: we obtain
$K(71)\ge 2\,603\,658\,750$, $K(70)\ge 1\,249\,778\,250$,
$K(69)\ge 627\,822\,180$ and $K(68)\ge 361\,275\,480$. In dimensions $46$ and
$47$ the same programme returns the published values $12\,309\,600$ and
$23\,766\,960$, both of them exactly, which validates the method.

The layered construction places the minimal shell on an equator of
$\R^{n}\oplus\R^{k}$ and adjoins lifted copies of part of it. Its level $t$ is a
free parameter, bounded below by $1/(2(1-\gamma))$, where $\gamma$ is the largest
inner product allowed between two minimal lines lifted above the same direction.
For $\Leech$ and $\Gsz$, $\gamma=\tfrac14$ and the level may be $\tfrac23$, at
which one set of lines serves a zero-sum triple of directions; for $\Pff$,
$\gamma=\tfrac13$ and only an antipodal pair is admissible. We obtain $K(25)\ge 197\,058$,
$K(27)\ge 200\,540$, $K(38)\ge 591\,612$ and improvements throughout
$49\le n\le 61$ and $73\le n\le 95$. Over $\Leech$ the same construction
reproduces the published values in dimensions $26$, $28$, $29$, $30$ and $31$.

The four remaining bounds are code-theoretic. Evaluated against the current code
tables, the chain of Edel--Rains--Sloane gives $K(39)\ge 756\,116$,
$K(62)\ge 71\,310\,732$ and $K(63)\ge 138\,419\,844$; at $n=96$ it overtakes
$\Gsz$, giving $K(96)\ge 12\,886\,999\,232$.

Every bound is certified in exact arithmetic by a public verification package;
no floating-point computation decides any of them.
\end{abstract}

\maketitle

\section{Introduction}\label{sec:intro}

The \emph{kissing number} $K(n)$ is the largest number of non-overlapping unit
balls that can touch a central unit ball in $\R^n$. Equivalently, it is the
largest size of a set of unit vectors in $\R^n$ whose pairwise inner products
are at most $\tfrac12$. Such a set is a \emph{kissing configuration}, or a
spherical code of minimum angle $60^\circ$.

Exact values are known only for $n=1,2,3,4,8,24$. The cases $n=3$
\cite{schutte-vanderwaerden1953}, $n=4$ \cite{musin2008}, $n=8$ and $n=24$
\cite{odlyzko-sloane1979,levenshtein1979} are hard theorems. In every other
dimension the known bounds leave an interval whose two ends move independently.
This paper concerns the lower end.

\subsection{The published tables}\label{ss:tables}

The state of the art is Cohn's \emph{Table of kissing number bounds}
\cite{cohn-table,cohn-data}, which supersedes the Nebe--Sloane table
\cite{nebe-sloane}. Their ranges fix the meaning of ``previously known'' below.
Cohn's table covers dimensions $1$--$48$ and $72$; the Nebe--Sloane table covers
$1$--$40$, $42$, $44$, $48$, $64$, $72$, $80$ and $128$. Above dimension $48$ the
two together have entries at $n=64$, $72$, $80$ and $128$ and at no other dimension.

In dimensions $49$--$63$ the best published value therefore comes from the
orthogonal direct sum of the minimal shell of $\Pff$ with a kissing configuration
of $\R^{n-48}$ and equals $52\,416\,000+K(n-48)$, where $K(n-48)$ is Cohn's entry in
dimension $n-48$, the argument running from $1$ to $15$. In dimensions $65$--$71$ it is
$331\,737\,984$, inherited by monotonicity of $K$ from the dimension-$64$ entry
of Edel--Rains--Sloane \cite{edel-rains-sloane1998}. In dimensions $73$--$96$ the
same construction over $\Gsz$ gives $6\,218\,175\,600+K(n-72)$, which at $n=96$ is
$6\,218\,175\,600+196\,560$, the second summand being the Leech lattice; here
$K(n-72)$ is again Cohn's entry, its argument running from $1$ to $24$. Forty-three of
the forty-seven improvements below lie in these ranges.

One dimension inside those three ranges is tabulated: at $n=80$ Nebe--Sloane record
$\Gsz\perp E_8$, whose kissing number $6\,218\,175\,600+240$ is the direct sum
above. The rule and the table agree there, and that is the only place they meet.

A table entry is not the only thing a bound above dimension $48$ has to beat. The
construction of Edel--Rains--Sloane is generic in $n$, and evaluated in the dimension
itself it can exceed the value inherited from a smaller one: in dimensions $65$--$71$
it does, by a little, giving $331\,771\,800$ at $n=68$ (Remark~\ref{rem:ers}). The
``previously'' column of Table~\ref{tab:all} is the tabular floor defined here;
\S\ref{ss:exposure} measures every bound of this paper against the constructive one,
which is the stronger test and the one on which the dimension-$68$ bound is tightest.

\subsection{Cross-sections and layers}

Both constructions start from the minimal shell of one of the three extremal
even unimodular lattices
\[
\begin{array}{lll}
  \Leech & n=24,\ \mu=4, & N=196\,560,\\
  \Pff   & n=48,\ \mu=6, & N=52\,416\,000,\\
  \Gsz   & n=72,\ \mu=8, & N=6\,218\,175\,600,
\end{array}
\]
where $\mu$ is the minimum and $N$ the number of minimal vectors. The minimal
shell of each, rescaled to the unit sphere, is a kissing configuration; $\Leech$
attains $K(24)=196\,560$.

The first construction takes cross-sections. The minimal vectors orthogonal to
$k$ fixed minimal vectors span a subspace of dimension at most $n-k$ and have
pairwise inner products at most $\tfrac12$; they therefore form a kissing
configuration in $\R^{n-k}$. The problem is to count them, for which Venkov's
theorem provides the constraints (\S\ref{sec:cross}).

The second construction places the shell on an equator of $\R^{n}\oplus\R^{k}$
and adjoins copies of part of it, lifted out of the equatorial subspace above a
small set of directions in $\R^{k}$ (\S\ref{sec:cap}). Every published record in
dimensions $25$ through $31$ belongs to this family.

A cross-section loses points and gives a bound below dimension $n$; the layered
construction gains points and gives one above it. The same quantity limits both:
how many minimal lines of $L$ can be pairwise at small inner product.

\subsection{Results}

Table~\ref{tab:summary} groups the improvements by construction.
Table~\ref{tab:all} in \S\ref{sec:results} gives one row per dimension.

\begin{table}[ht]
\caption{The forty-seven improvements, grouped by construction; $\#$ is the number
of dimensions in a row. The gain is $K_{\mathrm{here}}-K_{\mathrm{prev}}$ and the
factor their ratio, computed dimension by dimension, with both extremes
attained at an end of the dimension range, to the two decimals shown. ``Previously'' is \cite{cohn-table}
where that table has an entry and otherwise the value of \S\ref{ss:tables}.
Dimensions $44$ to $47$ are absent, for the reasons given with
Table~\ref{tab:all}, which gives every dimension.}
\label{tab:summary}
\small\setlength{\tabcolsep}{4.5pt}
\begin{tabular}{llrrrl}
\toprule
construction & dimensions & \# & largest gain & factor & where\\
\midrule
$\Gsz$ cross-sections & $68$--$71$ & $4$ & $+2\,271\,920\,766$ & $1.09$--$7.85$ & \S\ref{ss:g72}\\
layers over $\Leech$ & $25$, $27$, $38$ & $3$ & $+24\,960$ & $1.00$--$1.04$ & \S\ref{ss:d25}--\ref{ss:d38}\\
layers over $\Pff$ & $49$--$61$ & $13$ & $+7\,481\,540$ & $1.00$--$1.14$ & \S\ref{ss:big}\\
the \textsc{ers} chain & $62$, $63$ & $2$ & $+86\,001\,280$ & $1.36$--$2.64$ & \S\ref{ss:d6263}\\
the \textsc{ers} chain at $96$ & $96$ & $1$ & $+6\,668\,627\,072$ & $2.07$--$2.07$ & \S\ref{ss:d96}\\
layers over $\Gsz$ & $73$--$95$ & $23$ & $+254\,796\,296$ & $1.00$--$1.04$ & \S\ref{ss:big}\\
constant-weight codes & $39$ & $1$ & $+128$ & $1.00$--$1.00$ & \S\ref{ss:d39}\\
\midrule
total & & $47$ & & &\\
\bottomrule
\end{tabular}
\end{table}

In Table~\ref{tab:summary} the factor is largest where the published table has no
entry at all. Three rows exceed $2$, and in none of them is the gain a new
construction: the cross-sections count a classical configuration, and dimensions
$62$, $63$ and $96$ evaluate a construction of $1998$ against the code tables of
$2026$. The rows
with factors near $1$ span the most dimensions, $40$ of the $47$, though the two
largest absolute gains lie in rows whose factor exceeds $2$.

\subsection{Attribution}\label{ss:new}

\begin{sloppypar}
\emph{Dimensions $46$ and $47$.} The value $23\,766\,960$ is due to
Boyvalenkov--Cherkashin \cite{boyvalenkov-cherkashin2025}: it is equation~(3) of
the preprint \texttt{arXiv:2312.05121}, whose surrounding paragraph observes that
those $A_0$ vectors form a kissing configuration in $\R^{47}$.
The value $12\,309\,600$ is Table~3 of Sun--Wang \cite{sun-wang2026}
(\texttt{arXiv:2607.20359v3}), obtained by
enumerating the $52\,416\,000$ minimal vectors of $\Pff$; the degree-$3$ Siegel
theta tables of Ozeki \cite{ozeki2016} contain the ingredients a decade earlier.
We derive both in \S\ref{ss:p48} as a check on the cross-section programme, but
claim neither. Cohn's table \cite{cohn-table} records
$5\,318\,060$ in dimension $46$ and $9\,741\,412$ in dimension $47$.
\end{sloppypar}

\emph{The forty-seven listed below.} We have found no earlier claim in dimensions
$25$, $27$, $38$, $39$, $49$--$63$, $68$--$71$ or $73$--$96$. The one-point
distribution of the minimal shell of $\Gsz$ follows from Venkov's theorem in a
few lines; the bounds in dimensions $70$ and $71$ may therefore already be known
to specialists, though we have found no record of them. Dimensions $62$, $63$ and $96$
carry no new mathematics either. They evaluate the construction of
\cite{edel-rains-sloane1998} at three lengths where we find no record of its having
been evaluated: its seven worked examples are dimensions $32$, $36$, $40$, $44$,
$64$, $80$ and $128$. All three improve whenever the code tables do. The range
``$32$ to $128$'' of that paper's abstract is the span of those seven examples; it
tabulates no value by dimension, and dimension $96$ does not appear in it. Its own
totals depend on the code tables in the same way, and it says in closing that it
does not expect them to stand as those tables improve.

\subsection{Organisation}

\S\ref{sec:related} places this work in the literature. \S\ref{sec:prelim} fixes
notation and states the two classical inputs. \S\ref{sec:cross} develops the
cross-section method and proves the bounds in dimensions $68$ through $71$.
\S\ref{sec:cap} analyses the layered construction and proves the remaining
bounds, apart from dimensions $39$, $62$, $63$ and $96$, which are code-theoretic and
are treated in \S\ref{sec:ers}. \S\ref{sec:results} collects the full table.
\S\ref{sec:remarks} records the constructions we find to be exhausted, the margin
of the bounds above dimension $48$, one open question and the verification.

A public software package \cite{repo} verifies every bound below, to the
standards set out in \S\ref{ss:verif}; no floating-point computation decides any
of them.

\section{Related work}\label{sec:related}

\subsection{Upper bounds}\label{ss:rw-upper}

The results below concern lower bounds. The upper end is dominated by linear and
semidefinite programming. The method of Delsarte \cite{delsarte1973}, in the spherical
form of Delsarte, Goethals and Seidel \cite{delsarte-goethals-seidel1977}, yields the
classical bounds of Odlyzko and Sloane \cite{odlyzko-sloane1979} and of Levenshtein
\cite{levenshtein1979}, which determine $K(8)$ and $K(24)$. Semidefinite relaxations
strengthen it. Bachoc and Vallentin \cite{bachoc-vallentin2008} introduced the
three-point bound; Mittelmann and Vallentin \cite{mittelmann-vallentin2010} computed it
to high accuracy; Machado and de Oliveira Filho \cite{machado-oliveira2018} reduced it by
polynomial symmetry. The clustered low-rank solvers of Leijenhorst and de Laat
\cite{leijenhorst-delaat2024} produce most of the upper-bound column of
\cite{cohn-table} from dimension $10$ upwards and establish the optimality of the
$D_4$ root system \cite{delaat-leijenhorst-demuinckkeizer2024}. The gap remains wide in
the dimensions treated here: in dimension $47$ the best known construction gives
$23\,766\,960$ (\S\ref{ss:p48}) against an upper bound of $621\,658\,419$.

\subsection{Sphere packing}\label{ss:rw-packing}

The kissing number problem is the local form of the sphere packing problem. The two
have shared techniques since Cohn and Elkies \cite{cohn-elkies2003} introduced the
linear programming bound for packing density. Viazovska \cite{viazovska2017} settled
the packing problem in dimension $8$, and Cohn, Kumar, Miller, Radchenko and
Viazovska \cite{cohn-kumar-miller-radchenko-viazovska2017} in dimension $24$; the same
authors later proved universal optimality of $E_8$ and $\Leech$
\cite{cohn-kumar-miller-radchenko-viazovska2022}. Each of these proofs exhibits a modular
form making the linear programme tight. Dimensions $8$ and $24$ are exactly the
dimensions in which the kissing number was already known. No such form is available
in the dimensions considered here.

\subsection{Constructions from codes}\label{ss:rw-codes}

The entries of \cite{cohn-table} for $32\le n\le47$ come from the
Edel--Rains--Sloane construction \cite{edel-rains-sloane1998}, with the single
exception of dimension $43$. That construction assembles a $60^\circ$ code from a
chain of binary constant-weight codes and is therefore only as strong as the
code tables it is evaluated with \cite{brouwer}. Write $A(n,d)$ for the largest
size of a binary code of length $n$ with minimum Hamming distance $d$, and
$A(n,d,w)$ for the largest size of such a code all of whose words have weight
$w$. Lower down, Zinoviev and Ericson
\cite{zinoviev-ericson1999} improved a range of small dimensions by an extension scheme
built on Steiner systems. Of their values only dimension $13$ still stands; the work
of \S\ref{ss:rw-algebra} and \S\ref{ss:rw-search} has overtaken the rest.
Improvements to the underlying code tables propagate directly. Echols
\cite{echols2026} reports $124$ new constant-weight codes together with the kissing
numbers they give in dimensions $32$, $33$, $34$ and $37$; \S\ref{ss:d39} is a
further instance. Sun and Wang \cite{sun-wang2026} obtain dimensions $39$, $43$ and $45$ by an
antipode construction applied to cross-sections of $\Pff$. The same work
recomputes the kissing numbers of twelve Conway--Sloane cross-sections, including
the two recorded in \S\ref{ss:new}.

\subsection{Algebraic constructions}\label{ss:rw-algebra}

Ganzhinov \cite{ganzhinov2025} constructs highly symmetric line systems from twisted
spherical functions of finite groups and obtains kissing configurations of
$510$, $592$ and $1932$ points in dimensions $10$, $11$ and $14$. The
dimension-$14$ configuration provides the directions used in \S\ref{ss:d38}. Cohn
and Li \cite{cohn-li2024} improve dimensions $17$ through $21$ by reversing the
parity of the sign count in a Steiner-system construction, which admits further
vectors indexed by a subcode of a punctured Golay code. Ho \cite{ho2026} improves
their dimension-$19$ value to $11\,948$ by exhibiting a nonlinear subcode of that
kind. \S\ref{ss:obstructions} analyses both. In dimensions $25$ through $31$ the
predecessor is Kallal, Kan and Wang \cite{kallal-kan-wang2017}, whose configurations
belong to the family of \S\ref{sec:cap} and were superseded by those of Ma et
al.\ (\S\ref{ss:rw-search}).

\subsection{Numerical and machine-assisted search}\label{ss:rw-search}

A separate line of work obtains configurations by numerical optimisation, and
several of the values improved on below come from it.

Ma et al.\ \cite{ma-etal2025} apply game-theoretic reinforcement learning and hold
dimensions $25$ through $31$. Their configurations are the input to \S\ref{ss:d25}
and \S\ref{ss:d27}, which construct no new spherical code. Takhanov, Assylbekov and Yun
\cite{takhanov-assylbekov-yun2026} obtain $K(12)\ge841$ by energy minimisation initialised from
an analysis of the structure common to the known $840$-point arrangements.

Evolutionary search over programs driven by large language models has improved a
single dimension. AlphaEvolve \cite{novikov-etal2025} raised $K(11)$ from $592$ to
$593$; agents on the EinsteinArena platform \cite{bianchi-kwon-pappu-zou2026} later
raised it to $604$. Georgiev, G\'omez-Serrano, Tao and Wagner
\cite{georgiev-gomezserrano-tao-wagner2025} applied AlphaEvolve to $67$ problems
and examined the kissing number in dimensions $1$ through $11$. They recover the
known lower bound in each dimension and report $593$ in dimension $11$, adding no
further improvement; dimension $12$ was not attempted. GigaEvo
\cite{khrulkov-etal2025} attempted dimension $12$ with the aim of improving on the
lower bound $840$ and reached $840$.

The same means have been applied to a related problem. In
\cite{kravatskiy-khrulkov-oseledets2026}, by the present author with Khrulkov and
Oseledets, the object of study is the spherical code, or Tammes, problem: for fixed
$N$ and $d$, arrange $N$ points on $S^{d-1}$ so as to minimise the maximum pairwise
inner product. That work samples $90$ configurations at random from the catalogue of best known
spherical codes maintained by Cohn \cite{cohn-spherical}, in dimensions $8$
through $16$. It lowers the maximum inner product for the majority of them, by up
to $2.4\%$ in relative terms, and reports no kissing number result; the question
was not asked there.

\subsection{Scope of the present work}\label{ss:rw-position}

Numerical search has been effective where a configuration is small enough to be
stored as a list of coordinates and refined numerically; recent work of this kind
has concentrated on dimension $11$. The configurations treated here contain between
$1.98\times10^5$ and $6.4\times10^9$ points. A lattice together with a few integer
parameters specifies each of them. They lie beyond the reach of a
coordinate search. We obtain them from the structure of three extremal lattices.

Forty-three of the forty-seven dimensions below lie in ranges in which the published
tables have one entry between them, at $n=80$ (\S\ref{ss:tables}). The improvements there measure the extent
to which those ranges have not previously been examined.

The methods of this paper yield nothing in dimensions $9$ through $23$. Every
published record in dimensions $9$ through $19$ is maximal, in the sense that no
further point may be added to it, as are the cross-sections $\Lambda_{21}$,
$\Lambda_{22}$ and $\Lambda_{23}$; and the Cohn--Li construction is at its exact
ceiling in dimensions $18$ through $21$ (\S\ref{ss:obstructions}).


\section{Preliminaries}\label{sec:prelim}

Throughout, $L$ is an even unimodular lattice of dimension $n$ with minimum
$\mu$, and $\Lambda\subset L$ is its set of minimal vectors, or \emph{minimal
shell}, with $N=|\Lambda|$. For $v,w\in\R^n$ we write $\ip{v}{w}$ for the
standard inner product and $\hat v=v/|v|$.

We write $\Lambda_k$ for the laminated lattice of dimension $k$, following
\cite[Ch.~6]{conway-sloane1999}, so that $\Leech=\Lambda_{24}$; the unsubscripted
$\Lambda$ always denotes the minimal shell of $L$.

Every result below is a lower bound for $K$. Where the symbol $K(k)$ occurs
inside a construction, the inequality proved holds with any lower bound for
$K(k)$ substituted. The numerical values quoted use the best known one.

\begin{lemma}[the shell is a kissing configuration]\label{lem:shell}
$\{\hat v: v\in\Lambda\}$ is a kissing configuration in $\R^n$; hence
$K(n)\ge N$.
\end{lemma}

\begin{proof}
For distinct $v,w\in\Lambda$ the difference $v-w$ is a non-zero lattice vector,
so $|v-w|^2=2\mu-2\ip{v}{w}\ge\mu$, giving $\ip{v}{w}\le\mu/2$ and hence
$\ip{\hat v}{\hat w}\le\tfrac12$.
\end{proof}

$L$ is \emph{extremal} if its minimum is the largest permitted by the theory of
modular forms, namely $\mu=2\lfloor n/24\rfloor+2$.

\begin{proposition}[extremal theta series]\label{prop:theta}
The theta series of an extremal even unimodular lattice of dimension $n$ is the
unique modular form of weight $n/2$ for $\mathrm{SL}_2(\Z)$ with constant term
$1$ whose coefficients of $q^1,\dots,q^{\lfloor n/24\rfloor}$ vanish. In
particular $(\mu,N)$ depends only on $n$:
\[
\begin{array}{lcccc}
 n   & 8   & 24      & 48           & 72 \\
 \mu & 2   & 4       & 6            & 8 \\
 N   & 240 & 196\,560 & 52\,416\,000 & 6\,218\,175\,600.
\end{array}
\]
\end{proposition}

Four extremal lattices are known in dimension $48$: $P_{48p}$ and $P_{48q}$
\cite[Ch.~7]{conway-sloane1999}, together with $P_{48m}$ and $P_{48n}$. Nebe
constructed the fourth of them \cite{nebe2014}; the catalogue
\cite{nebe-sloane} records all four. One extremal lattice is known in dimension
$72$, the lattice $\Gsz$ of Nebe \cite{nebe2012}. Nothing in \S\ref{sec:cross}
distinguishes the four, the bounds there being read off the Gram matrix and the
one-point distribution, and we write $\Pff$ for any of them; the searches on
explicit coordinates in \S\ref{ss:big} are carried out in $P_{48p}$.

\begin{proposition}[Venkov \cite{venkov1984}]\label{prop:venkov}
If $L$ is extremal even unimodular of dimension $n\equiv 0\pmod{24}$, then
$\{\hat v: v\in\Lambda\}$ is a spherical $11$-design: every polynomial of
degree at most $11$ has the same average over $\Lambda$ as over the sphere.
For $n=8$ the shell is a $7$-design.
\end{proposition}

Integrality bounds the inner products in the shell.

\begin{lemma}\label{lem:ips}
For $v,w\in\Lambda$ the inner product $\ip{v}{w}$ is an integer, and
$|\ip{v}{w}|\le\mu/2$ unless $w=\pm v$, in which case $\ip{v}{w}=\pm\mu$.
\end{lemma}

\begin{proof}
Integrality because $L$ is integral. Both $|v-w|^2=2\mu-2\ip{v}{w}$ and
$|v+w|^2=2\mu+2\ip{v}{w}$ are norms of lattice vectors, so each is either $0$ or
at least $\mu$. If neither vanishes, the two bounds give $|\ip{v}{w}|\le\mu/2$;
and $|v\mp w|^2=0$ exactly when $w=\pm v$, which gives $\ip{v}{w}=\pm\mu$.
\end{proof}

Thus for $\Pff$ the inner products lie in $\{0,\pm1,\pm2,\pm3\}\cup\{\pm6\}$, and
for $\Gsz$ in $\{0,\pm1,\pm2,\pm3,\pm4\}\cup\{\pm8\}$.

\section{Cross-sections of extremal lattices}\label{sec:cross}

\subsection{Cross-sections and their cells}

\begin{lemma}\label{lem:cross}
Let $u_1,\dots,u_k\in\Lambda$ be linearly independent and put
\[
  S=\{\,v\in\Lambda : \ip{u_i}{v}=0 \text{ for } i=1,\dots,k\,\}.
\]
Then $S$ lies in a subspace of dimension $n-k$ and $K(n-k)\ge|S|$. We call $S$
the \emph{cross-section} of $\Lambda$ by $u_1,\dots,u_k$.
\end{lemma}

\begin{proof}
$S$ lies in $\mathrm{span}(u_1,\dots,u_k)^\perp$. All its members have norm
$\mu$; Lemma~\ref{lem:shell} applies verbatim.
\end{proof}

It remains to count $S$. By Lemma~\ref{lem:ips} the vector
$(\ip{u_1}{v},\dots,\ip{u_k}{v})$ attached to $v\in\Lambda$ lies in
$\mathcal I^k$, where $\mathcal I=\{0,\pm1,\dots,\pm\mu/2\}\cup\{\pm\mu\}$. We
call an element $c=(c_1,\dots,c_k)$ of $\mathcal I^k$ a \emph{cell} and write
\[
  n[c]\;=\;n[c_1,\dots,c_k]
  \;=\;\#\{\,v\in\Lambda : \ip{u_i}{v}=c_i \text{ for } i=1,\dots,k\,\}
\]
for the number of minimal vectors it holds, the \emph{count} of $c$. The notation
leaves $u_1,\dots,u_k$ implicit, one cross-section being under discussion at a
time; at $k=1$ its argument is a single integer. Thus $|S|=n[0,\dots,0]$, the
counts of the $|\mathcal I|^k$ cells summing to $N$. The family
$\bigl(n[c]\bigr)_{c\in\mathcal I^k}$ is the \emph{$k$-point
distribution} of $u_1,\dots,u_k$.

\subsection{Moments from the design property}

\begin{lemma}\label{lem:moments}
For every exponent vector $p=(p_1,\dots,p_k)$ of non-negative integers with
$\sum p_i\le 11$ and $\sum p_i$ even,
\begin{equation}\label{eq:moment}
  \sum_{c} n[c]\prod_{i=1}^{k}\Bigl(\frac{c_i}{\mu}\Bigr)^{p_i}
  \;=\;
  \frac{N\cdot \mathbb{E}\bigl[\prod_i \ip{\hat u_i}{\zeta}^{p_i}\bigr]}
       {n(n+2)\cdots(n+\textstyle\sum p_i-2)},
\end{equation}
where $\zeta$ is a standard Gaussian vector in $\R^n$ and the numerator is
evaluated from the normalised Gram matrix $\bigl(\ip{\hat u_i}{\hat u_j}\bigr)$
by Wick's theorem.
\end{lemma}

\begin{proof}
The monomial $x\mapsto\prod_i\ip{\hat u_i}{x}^{p_i}$ has degree $\sum p_i\le11$,
so by Proposition~\ref{prop:venkov} its sum over the rescaled shell equals $N$
times its average over the sphere. Writing $x=\zeta/|\zeta|$ and using the
independence of $|\zeta|$ and $\zeta/|\zeta|$ gives the normalisation
$\mathbb{E}|\zeta|^{2m}=n(n+2)\cdots(n+2m-2)$. The Gaussian moment is a sum over
perfect matchings, each pair contributing a Gram entry. Odd total degree
vanishes because $\Lambda=-\Lambda$. Finally
$\ip{\hat u_i}{\hat v}=\ip{u_i}{v}/\mu$ since all vectors have norm $\mu$.
\end{proof}

\begin{remark}[mixed exponents]\label{rem:mixed}
Imposing \eqref{eq:moment} only for those $p$ with every $p_i$ even gives a valid
but weaker programme, whose value depends on the presentation of the Gram matrix.
At $k=3$ the full family has $161$ equations and the even-exponent family $56$.
Write a $3\times3$ Gram matrix of minimal vectors by its off-diagonal entries.
Write $3A_3$ for three times the Gram matrix of the root lattice $A_3$. The
matrix $G=(3,-3,0)$ satisfies $U^{\mathsf T}GU=3A_3$ for an integer matrix $U$
with $\det U=-1$; thus $G$ and $3A_3$ describe the same configuration. With even
exponents alone the two presentations give $6\,462\,480$ and $5\,194\,969$; with
the full family both give $6\,687\,320$.
\end{remark}

\subsection{Constraints from the lattice structure}

Proposition~\ref{prop:venkov} uses only that $\Lambda$ is a spherical
$11$-design. The next two lemmas sharpen the bound by using that $\Lambda$ is the
minimal shell of a lattice.

\begin{lemma}[translations]\label{lem:trans}
Suppose $w=\sum a_iu_i$ is a lattice vector with $|w|^2=\mu$, and let
$\tau_l=\ip{u_l}{w}$. If $\ip{w}{v}=\mu/2$ for $v\in\Lambda$, then
$w-v\in\Lambda$ and $v\mapsto w-v$ is an involution. Consequently
\[
  n[c]=n[c'], \qquad c'_l=\tau_l-c_l ,
\]
for every cell $c$ with $\sum_i a_ic_i=\mu/2$. If $c'\notin\mathcal I^k$ the constraint
becomes $n[c]=0$.
\end{lemma}

\begin{proof}
$|w-v|^2=2\mu-2\ip{w}{v}=\mu$, so $w-v$ is minimal, with
$\ip{u_l}{w-v}=\tau_l-c_l$. The hypothesis $\ip{w}{v}=\sum a_ic_i=\mu/2$ depends
only on the cell; the involution therefore maps the cell $c$ onto the cell $c'$.
\end{proof}

\begin{lemma}[lattice range]\label{lem:range}
For every lattice vector $w=\sum a_iu_i$ in the span, with $|w|^2=\nu>0$, and
every $v\in\Lambda$,
\[
  |\ip{v}{w}|\le \nu/2 \qquad\text{unless } v=\pm w \text{ (possible only if }\nu=\mu).
\]
A cell violating this for some such $w$ is empty.
\end{lemma}

\begin{proof}
$|v\mp w|^2=\mu+\nu\mp2\ip{v}{w}$ is a lattice norm, hence $0$ or at least
$\mu$.
\end{proof}

Lemma~\ref{lem:range} is strictly stronger than testing positive
semidefiniteness of the Gram matrix of $(u_1,\dots,u_k,v)$: for a $60^\circ$
pair, $u_1-u_2$ is itself minimal, which forces $|c_1-c_2|\le\mu/2$ and excludes
ten cells that semidefiniteness permits. Those ten cells are the zeros of
\cite[Lemma 4.2]{ozeki2016}.

Two further exact constraints are available. The first is the antipodal
symmetry $n[c]=n[-c]$. The second is the set of one-point marginals, which are
Lemma~\ref{lem:moments} at $k=1$; \S\ref{ss:onepoint} solves them in closed form.

\subsection{The linear programme and its dual certificate}\label{ss:cert}

Call the linear programme of this subsection, built from $k$ fixed vectors, the
\emph{$k$-point programme}. Its equality constraints form the \emph{$k$-point
system}.

For a cell $c$ let $G(c)$ be the Gram matrix of $(u_1,\dots,u_k,v)$ for a vector
$v$ of norm $\mu$ with $\ip{u_i}{v}=c_i$ for every $i$. Its entries are
determined by $c$ and by the Gram matrix of $u_1,\dots,u_k$, whether or not such
a $v$ lies in $\Lambda$.

Collect the constraints as $Ax=b$, $0\le x\le h$, with $x$, $h$ and $s$ below all
indexed by cells. Set $h_c=1$ when $\det G(c)=0$ and $h_c=+\infty$ otherwise.
The first case is legitimate: a vector $v\in\Lambda$ counted by such a cell lies
in $\mathrm{span}(u_1,\dots,u_k)$ and is determined there by $c$; at most one
such $v$ therefore exists. Let $e_0$ be the indicator of the all-zero cell. The
true distribution is feasible, whence
\[
  |S| \;\ge\; \min\{\, e_0^{\mathsf T}x : Ax=b,\ 0\le x\le h \,\}.
\]

\begin{lemma}[weak duality]\label{lem:dual}
For any rational vector $y$, put $s=e_0-A^{\mathsf T}y$. If $s_c\ge0$ for every
cell $c$ with $h_c=+\infty$, then
\[
  \min\{\,e_0^{\mathsf T}x : Ax=b,\ 0\le x\le h\,\}
  \;\ge\;
  b^{\mathsf T}y-\sum_{c:\,h_c<\infty}h_c\max(0,-s_c).
\]
\end{lemma}

\begin{proof}
For feasible $x$ we have
$e_0^{\mathsf T}x=y^{\mathsf T}Ax+s^{\mathsf T}x=b^{\mathsf T}y+s^{\mathsf T}x$,
and $s^{\mathsf T}x\ge-\sum_{h_c<\infty}h_c\max(0,-s_c)$ because
$0\le x_c\le h_c$.
\end{proof}

A floating-point solver is used only to propose a candidate $y$. That $y$ is
rounded to rationals and repaired into exact dual feasibility by decreasing the
coordinate belonging to the all-ones moment row, whose coefficients are all
positive. The right-hand side of Lemma~\ref{lem:dual} is then evaluated in exact
rational arithmetic and rounded to an integer downwards. An imprecise $y$ weakens
the bound and cannot invalidate it. Applying the same lemma to $-e_0$ certifies an
upper bound on $|S|$, rounded upwards. A certified minimum may therefore fall one
below the true optimum, and a certified maximum one above it.

\subsection{Realizability}\label{ss:realize}

Call a symmetric $k\times k$ integer matrix $G$ with every diagonal entry equal
to $\mu$ \emph{realizable in} $L$ if some $u_1,\dots,u_k\in\Lambda$ have Gram
matrix $G$. A bound computed for $G$ is vacuous unless $G$ is realizable in $L$.

\begin{lemma}\label{lem:realize}
Let $u_1,\dots,u_{k-1}\in\Lambda$ have Gram matrix $G'$, and let
$c\in\mathcal I^{k-1}$.
Suppose the minimum of $n[c]$ over the $(k-1)$-point programme is strictly
positive. Then some $v\in\Lambda$ satisfies $\ip{u_i}{v}=c_i$ for
$i=1,\dots,k-1$, and the matrix obtained by bordering $G'$ with the row $c$ and
the diagonal entry $\mu$ is realizable in $L$.
\end{lemma}

\begin{proof}
The true distribution is feasible for the $(k-1)$-point programme; its value
at the cell $c$ is therefore at least the minimum, hence positive. Take $u_k=v$.
\end{proof}

Lemma~\ref{lem:realize} is the weaker of the two routes to realizability. A sweep
over all Gram matrices produces larger apparent bounds whose bordering cells have
minimum $0$, about which the lemma says nothing; among three-point Gram matrices
in dimension $69$ it licenses only the pairwise orthogonal one, worth
$149\,415\,783$, below the value $331\,737\,984$ of \S\ref{ss:tables}.
Realizability is, however, a property of the lattice and not of the programme.
Where the lattice is available in coordinates a $k$-tuple with a prescribed Gram
matrix can simply be \emph{exhibited}, as it is for $\Gsz$: dimensions $68$ and
$69$ of Table~\ref{tab:all} are obtained this way.

Which Gram matrix to exhibit is then guided by a measurement. Across the
Gram matrices tried here the bound falls as $\det G$ rises --- at $k=3$ from
$527\,597\,644$ at $\det G=352$ to $149\,415\,783$ at $512$. At a fixed determinant the
presentation costs nothing: at $k=3$ the sixteen realizable Gram matrices of determinant
$256$ are sixteen bases of one scaled $A_3$, and all sixteen certify $627\,822\,180$. Taking the
determinant as low as it will go, $u_1,\dots,u_k$ generate a sublattice $M$ of rank $k$ with
$\min M\ge\mu$ and $\det M=\det G$, so Hermite's inequality gives
$\det M\ge(\mu/\gamma_k)^k$, where $\gamma_k$ is Hermite's constant in rank $k$,
with equality exactly for the densest rank-$k$ lattice. For $\Gsz$, where
$\mu=8$, that is $48$ at $k=2$, $256$ at $k=3$ and $1024$ at $k=4$, attained by
$A_2$, $A_3$ and $D_4$. Those are the three Gram matrices used in dimensions
$70$, $69$ and $68$; each is exhibited inside $\Gsz$. Dimension $69$ gains $100\,224\,536$ over
the Gram $(4,-2,0)$, whose determinant is $352$.

\subsection{The one-point distribution}\label{ss:onepoint}

At $k=1$ the system of Lemma~\ref{lem:moments} has $\mu/2+1$ unknowns and six
equations, one for each even degree up to $10$; the shell of $E_8$ is a
$7$-design, which gives four equations there. The system is over-determined; the
surplus equations serve as a consistency check.

\begin{proposition}\label{prop:onepoint}
For $u\in\Lambda$ and $c=0,\dots,\mu/2$ let $n[c]$ be the number of $v\in\Lambda$
with $\ip{u}{v}=c$. Then
\begin{center}
\begin{tabular}{lrrrrr}
\toprule
lattice & $n[0]$ & $n[1]$ & $n[2]$ & $n[3]$ & $n[4]$\\
\midrule
$E_8$    & $126$ & $56$ & & & \\
$\Leech$ & $93\,150$ & $47\,104$ & $4\,600$ & & \\
$\Pff$   & $23\,766\,960$ & $12\,608\,784$ & $1\,678\,887$ & $36\,848$ & \\
$\Gsz$   & $2\,603\,658\,750$ & $1\,512\,243\,200$ & $280\,928\,256$ & $13\,959\,168$ & $127\,800$\\
\bottomrule
\end{tabular}
\end{center}
with each non-zero $c$ attained by $n[c]$ vectors of each sign. The surplus
equations hold in every case.
\end{proposition}

The first two rows are the root system $E_7$ and the laminated lattice
$\Lambda_{23}$, whose kissing numbers are known independently. The $\Gsz$ row and
Lemma~\ref{lem:cross} give Theorem~\ref{thm:71}; the $\Pff$ row recovers a
published value (\S\ref{ss:p48}).

\subsection{Dimensions 46 and 47}\label{ss:p48}

Both values below are in the literature (\S\ref{ss:new}). We derive them here to
test the programme against a published value.

\begin{proposition}\label{thm:47}
The cross-section of $\Pff$ by one minimal vector has exactly $23\,766\,960$
elements, and the cross-section by two minimal vectors at $60^\circ$ has exactly
$12\,309\,600$. Hence $K(47)\ge 23\,766\,960$ and $K(46)\ge 12\,309\,600$. Both
values are known, the first from \cite{boyvalenkov-cherkashin2025} and the second
from \cite{sun-wang2026}.
\end{proposition}

\begin{proof}
For $k=1$, Proposition~\ref{prop:onepoint} gives $n[0]=23\,766\,960$ and the
system determines it uniquely; Lemma~\ref{lem:cross} finishes.

For $k=2$ take $\ip{u_1}{u_2}=3$, a $60^\circ$ pair. Such a pair exists: the
one-point programme has value $n[3]=36\,848>0$ at that cell, so
Lemma~\ref{lem:realize} applies. Running the programme of Lemmas~\ref{lem:moments}--\ref{lem:dual} on
the all-zero cell, first minimising and then maximising, returns the same
certified value $12\,309\,600$ in both directions. Since the true size lies
between the two, it equals $12\,309\,600$.
\end{proof}

For the orthogonal pair the programme gives an interval,
\[
  10\,659\,120 \;\le\; |S| \;\le\; 10\,697\,520,
\]
and no bound depending only on the Gram matrix can do better. In the notation of
Remark~\ref{rem:ozeki}, row $20$ of \cite[Table 3]{ozeki2016} divided by
$a((3,3,0))$ equals $458\,386\,320/43$, which is not an integer. That quotient is
an exact average over ordered orthogonal pairs, so the size varies from pair to
pair. The width of the interval reflects that variation.

\begin{remark}[modular forms]\label{rem:ozeki}
The dimension-$46$ value also follows from modular forms, with no linear
programming. Write $a(T)$ for the Siegel Fourier coefficient of the theta series
of $\Pff$ at the half-integral form $T$, which counts the tuples of lattice
vectors whose Gram matrix is $2T$; here $T$ is listed as its diagonal entries
followed by the off-diagonal entries of $2T$, so that minimal vectors give
diagonal $3$ and pairwise entries $\ip{u_i}{u_j}$. The degree-$3$ coefficient at
a $60^\circ$ pair together with a vector orthogonal to both, over the degree-$2$
coefficient at a $60^\circ$ pair, is
\[
  \frac{a\bigl((3,3,3,3,0,0)\bigr)}{a\bigl((3,3,3)\bigr)}
  = \frac{23\,775\,066\,324\,172\,800\,000}{1\,931\,424\,768\,000}
  = 12\,309\,600
\]
exactly, with no remainder; the degree-$3$ coefficients are
\cite[Table 3]{ozeki2016}. Sun--Wang \cite{sun-wang2026} obtain the same value by
direct enumeration.
\end{remark}

\subsection{Dimensions 68 to 71}\label{ss:g72}

\begin{theorem}\label{thm:71}
$K(71)\ge 2\,603\,658\,750$ and $K(70)\ge 1\,249\,778\,250$. The first is the
exact size of that cross-section of $\Gsz$.
\end{theorem}

\begin{proof}
Since $72\equiv0\pmod{24}$, the shell of $\Gsz$ is an $11$-design by
Proposition~\ref{prop:venkov}. For a fixed $u\in\Lambda$, Lemma~\ref{lem:ips}
allows
$\ip{u}{v}=c$ with $|c|\le4$, together with $c=\pm8$ for $v=\pm u$: five
unknowns against six moment equations. The system is over-determined, solves in
non-negative integers, and the surplus equation holds
(Proposition~\ref{prop:onepoint}). Then $\Gsz\cap u^\perp$ is a
$71$-dimensional lattice of minimum $8$ whose minimal vectors are exactly the
$n[0]$ vectors orthogonal to $u$; Lemma~\ref{lem:cross} gives the first
bound.

For $k=2$ take $\ip{u_1}{u_2}=4$. Such a pair exists because $n[4]=127\,800>0$.
The two-point programme certifies $1\,249\,778\,250$. Sweeping all five admissible inner
products $\ip{u_1}{u_2}=0,1,2,3,4$ gives, in that order, $1\,073\,551\,500$,
$1\,089\,616\,500$, $1\,117\,183\,200$, $1\,167\,565\,500$ and
$1\,249\,778\,250$; the $60^\circ$ pair is the best of them.
\end{proof}

Unlike dimension $46$, the minimum and maximum do not meet here: the interval is
$[1\,249\,778\,250,\,1\,250\,506\,441]$, of which only the lower end is claimed.

At $k=3$ and $k=4$ the moment system no longer forces the bound and the Gram
matrix has to be chosen; \S\ref{ss:realize} takes one at the Hermite floor. That
choice is guided by a measurement, but the bound does not rest on it: each
certificate is valid for the Gram matrix it is computed from, and that matrix is
exhibited.

\begin{theorem}\label{thm:6869}
$K(69)\ge 627\,822\,180$ and $K(68)\ge 361\,275\,480$.
\end{theorem}

\begin{proof}[Proof sketch]
For $k=3$ take the Gram matrix with diagonal $8$ and off-diagonal entries
$-4,-4,0$, that of $A_3$ scaled by $4$, of determinant $256$. For $k=4$ take that
of $D_4$ scaled by $4$, of determinant $1024$. By \S\ref{ss:realize} these are
the least determinants a triple and a quadruple of minimal vectors of $\Gsz$ can
have. Both are realizable by exhibition. Four minimal vectors of $\Gsz$ with the
second Gram matrix are written out in \cite{repo}; the rank-$4$ sublattice they
generate has $24$ minimal vectors, among which a triple with the first is
displayed. The three-point and four-point programmes then certify $627\,822\,180$
and $361\,275\,480$; Lemma~\ref{lem:cross} finishes.
\end{proof}

What a numerically proposed dual can be repaired into depends on the basis; the
programme's optimum does not. Of the $104$ quadruples of those $24$ minimal
vectors whose Gram matrix has determinant $1024$, ninety-eight yield a repaired
certificate, running from $358\,774\,871$ to $361\,275\,383$, a spread of $0.7\%$;
the solver does not finish the other six. Solving for the dual \emph{vertex}
instead removes it. Substituting out the
constraints that are identifications --- $n[c]=n[-c]$, and the translations that
read $n[c]=n[c']$, $n[c]=0$ or $n[c]=1$ --- is exact and takes the $k=4$
programme from $3473$ rows and $1681$ columns to $617$ and $520$; the active
system is then solved in exact rationals and $A^{\mathsf T}y\le c$ verified over
every column. The optimum is
$361\,275\,480$, and the numeric primal attains it, so no dual improves on it.

The chain ends at $k=4$. The certified values fall by factors $0.480$, $0.502$
and $0.575$ as $k$ runs from $1$ to $4$; reaching $331\,737\,984$ at $k=5$ would
need a factor $0.92$, and the Hermite floor there is $4096$. We have not computed
$k=5$, where the cell enumeration is $11^k$; the trend is evidence that it would
fall short, not a proof. Dimensions $65$ through $67$ keep their inherited values
here, and we claim nothing in them.

\begin{remark}[the Edel--Rains--Sloane value]\label{rem:ers}
No table has an entry in dimensions $65$--$71$, so the value to improve on is
$331\,737\,984$, inherited from dimension $64$ \cite{edel-rains-sloane1998}. The
chain of \S\ref{sec:ers} reproduces both of their large published totals:
\begin{align*}
  2^{28}+30828\cdot2048+10416\cdot16+128 &= 331\,737\,984,\\
  2^{30}+143780\cdot2048+20540\cdot16+160 &= 1\,368\,532\,064.
\end{align*}
An inherited value is not the value the construction takes in the dimension that
inherits it. The difference matters in dimension $68$, where the margin is
smallest. There the chain is $(64,16,4,1)$, not $(68,\dots)$: the level-zero sign code
has length $64$, the best available at any length up to $68$, while the support codes
spend all $68$ coordinates. It gives $331\,771\,800$ from
$A(68,16,16)\ge A(64,16,16)\ge30\,828$ and
$A(68,4,4)=\binom{68}{3}/4=12\,529$. The second of these is exact, since
$68\equiv2\pmod6$ admits a Steiner system $S(3,4,68)$. Theorem~\ref{thm:6869}
exceeds that by a factor $1.089$, the narrowest margin over the chain of any bound
here not derived from it; \S\ref{ss:exposure} measures the others. Evaluated at $n=70$ and $n=71$, granting it the dimension-$80$
constant-weight code sizes it has no claim to at those lengths, the chain gives
at most $563\,125\,646$, which Theorem~\ref{thm:71} exceeds by factors $2.2$ and
$4.6$.
\end{remark}

\subsection{Validation against known cross-sections}\label{ss:validate}

An invalid constraint inflates the minimum of the programme and makes the lower
bound false. We therefore ran the programme on every cross-section whose size is
known independently and checked that the known size lies inside the interval the
programme produces. Table~\ref{tab:validate} reports the outcome.

\begin{table}[ht]
\caption{Every independently known cross-section, against the interval the
method produces for it. Sources of truth: classical values for the laminated
lattices; direct enumeration from coordinates; and the published computations
of \cite{boyvalenkov-cherkashin2025,ozeki2016,sun-wang2026}. The two orthogonal
Gram matrices over $\Leech$ admit embeddings with different counts: at
$k=3$ both are given, at $k=4$ ($\dagger$) one of several.}
\label{tab:validate}
\begin{tabular}{llrrr}
\toprule
lattice & configuration & minimum & \textbf{truth} & maximum \\
\midrule
$E_8$    & $k=1$                   & $126$ & $\mathbf{126}$ & $126$ \\
$E_8$    & $k=2$ orthogonal        & $60$ & $\mathbf{60}$ & $60$ \\
$E_8$    & $k=2$ at $60^\circ$     & $72$ & $\mathbf{72}$ & $72$ \\
$\Leech$ & $k=1$                   & $93\,150$ & $\mathbf{93\,150}$ & $93\,150$ \\
$\Leech$ & $k=2$ orthogonal        & $43\,164$ & $\mathbf{43\,164}$ & $43\,164$ \\
$\Leech$ & $k=2$ at $\cos\tfrac14$ & $44\,550$ & $\mathbf{44\,550}$ & $44\,550$ \\
$\Leech$ & $k=2$ at $60^\circ$     & $49\,896$ & $\mathbf{49\,896}$ & $49\,896$ \\
$\Leech$ & $k=3$ orthogonal        & $19\,530$ & $\mathbf{19\,530},\,\mathbf{19\,962}$ & $19\,962$ \\
$\Leech$ & $k=3$, $2A_3$           & $27\,720$ & $\mathbf{27\,720}$ & $27\,720$ \\
$\Leech$ & $k=4$ orthogonal        & $7\,464$ & $\mathbf{8\,616}^{\dagger}$ & $17\,400$ \\
$\Leech$ & $k=4$, $2A_4$           & $15\,540$ & $\mathbf{15\,540}$ & $15\,540$ \\
$\Leech$ & $k=4$, $2D_4$           & $17\,400$ & $\mathbf{17\,400}$ & $17\,400$ \\
$\Pff$   & $k=1$                   & $23\,766\,960$ & $\mathbf{23\,766\,960}$ & $23\,766\,960$ \\
$\Pff$   & $k=2$ at $60^\circ$     & $12\,309\,600$ & $\mathbf{12\,309\,600}$ & $12\,309\,600$ \\
$\Pff$   & $k=3$, $3A_3$           & $6\,687\,320$ & $\mathbf{6\,687\,648}$ & $6\,702\,736$ \\
$\Pff$   & $k=3$, $3D_3$           & $6\,687\,320$ & $\mathbf{6\,702\,080}$ & $6\,702\,736$ \\
$\Pff$   & $k=4$, $3D_4$           & $4\,147\,967$ & $\mathbf{4\,153\,600}$ & $4\,165\,632$ \\
$\Gsz$   & $k=1$                   & $2\,603\,658\,750$ & $\mathbf{2\,603\,658\,750}$ & $2\,603\,658\,750$ \\
\bottomrule
\end{tabular}
\end{table}

Every known size lies inside its interval.

Dimension $68$ is a cross-section at $k=4$, and on the Leech lattice the truth is
available there. The programme is tight at the two Gram matrices of least
determinant: for $2D_4$, which attains the Hermite floor $64$, and for $2A_4$, of
determinant $80$, the minimum and the maximum coincide with the count. The
orthogonal quadruple, of determinant $256$ against $D_4$'s $64$,
shows that this is a property of the good Gram matrices rather than of the
programme; at that Gram matrix no programme reading the Gram matrix alone could
be exact, the count not being a function of it.

The size of a cross-section is not a function of the Gram matrix. The rows $3A_3$
and $3D_3$ carry two bases of one lattice --- $A_3$ and $D_3$ are the same abstract
lattice --- so the programme returns one interval for both, yet their sizes differ:
$6\,687\,648$ and $6\,702\,080$, at very nearly its two ends.
Over $\Leech$ the effect can be exhibited outright: in the
coordinates in which a minimal vector has shape $(\pm4^2,0^{22})$, the triples
\[
  4e_1+4e_2,\quad 4e_3+4e_4,\quad 4e_5+4e_6
  \qquad\text{and}\qquad
  4e_1+4e_2,\quad 4e_1-4e_2,\quad 4e_3+4e_4
\]
are both pairwise orthogonal with every norm $\mu$, so they share the Gram matrix
$\mu I_3$, and their cross-sections have $19\,530$ and $19\,962$ elements ---
exactly the two ends of the certified interval. Orthogonal quadruples behave the
same way, realising $8\,616$, $8\,760$ and $9\,192$ among others. A programme
reading the Gram matrix alone can therefore only bracket such a row, and the
truth column records one embedding in each of these two.

The rows with $k=3$ also recover a theorem of Ozeki. The three-point system for a
$60^\circ$ triple has affine solution space of dimension $1$, along
$n[0,0,0]+164\,n[3,3,3]=6\,732\,912$. Ozeki's range $184\le n[3,3,3]\le278$
\cite[Theorem 6.3]{ozeki2016} therefore corresponds to
$6\,687\,320\le n[0,0,0]\le6\,702\,736$, which is exactly the interval certified
in Table~\ref{tab:validate}; the outward rounding of \S\ref{ss:cert} costs
nothing at either end. Ozeki obtains the range from modular forms; the programme
obtains it with an exact rational certificate.

\section{The layered construction}\label{sec:cap}

\subsection{The construction}

Fix a lattice $L$ as in \S\ref{sec:prelim} and an auxiliary dimension $k$, and
work in $\R^{n}\oplus\R^{k}$. Choose a level $t\in(0,1)$ and write
$a=\sqrt t$, $b=\sqrt{1-t}$.

A minimal \emph{line} of $L$ is an antipodal pair $\pm u$ with $u\in\Lambda$. Set
\begin{equation}\label{eq:gamma}
  \gamma=\frac{\lfloor \mu/3\rfloor}{\mu},
\end{equation}
the largest inner-product threshold available in $L$ that does not exceed
$\tfrac13$. Thus $\gamma=\tfrac14$ for $\Leech$ ($|\ip{u}{u'}|\le1$ of $4$) and
for $\Gsz$ ($\le2$ of $8$), while $\gamma=\tfrac13$ for $\Pff$ ($\le2$ of $6$).
Let $\Gamma_\gamma(L)$ be the graph on the minimal lines of $L$ in which two
lines are adjacent when the cosine of the angle between them exceeds $\gamma$ in
absolute value.

Let $C_1,\dots,C_r$ be pairwise disjoint independent sets of $\Gamma_\gamma(L)$,
and let $Z_i\subset S^{k-1}$ be a set of unit \emph{directions} attached to
$C_i$. The configuration is
\begin{equation}\label{eq:families}
\begin{array}{lll}
  \text{equator} & (\hat v,\,0) & v\in\Lambda,\ v\ne\pm u \text{ for any } u\in\textstyle\bigcup_i C_i,\\[2pt]
  \text{caps}    & (s\,a\,\hat u,\; b\,z) & s\in\{\pm1\},\ u\in C_i,\ z\in Z_i,\\[2pt]
  \text{poles}   & (0,\,w) & w\in W\subset S^{k-1}.
\end{array}
\end{equation}
Every point is a unit vector. We write $Z=\bigcup_i Z_i$ for the set of all
directions. Counting,
with $|C_i|$ the number of lines in $C_i$,
\begin{equation}\label{eq:count}
  \#\;=\;N-2\sum_i|C_i|\;+\;2\sum_i|C_i|\,|Z_i|\;+\;|W|
  \;=\;N+2\sum_i|C_i|\bigl(|Z_i|-1\bigr)+|W| .
\end{equation}

\begin{proposition}\label{prop:cap}
The configuration \eqref{eq:families} is a kissing configuration in
$\R^{n+k}$ if and only if
\begin{align}
  t &\le 1, \label{eq:c2}\\
  t\gamma+(1-t) &\le \tfrac12, \label{eq:c3}\\
  \ip{z}{z'} &\le \frac{\tfrac12-t}{1-t}
     \quad\text{for distinct } z,z'\in Z_i, \label{eq:c4}\\
  \ip{z}{z'} &\le \tfrac12
     \quad\text{for } z\in Z_i,\ z'\in Z_j,\ i\ne j, \label{eq:c5}\\
  b\,\ip{z}{w} &\le \tfrac12
     \quad\text{for } z\in Z,\ w\in W, \label{eq:c6}\\
  \ip{w}{w'} &\le \tfrac12
     \quad\text{for distinct } w,w'\in W. \label{eq:c7}
\end{align}
\end{proposition}

\begin{proof}
There are nine kinds of pair. Five are immediate. Equator against equator is
Lemma~\ref{lem:shell} and equator against pole is $0$. Equator against cap gives
$s\,a\ip{\hat v}{\hat u}$; since $v\ne\pm u$ forces
$|\ip{\hat v}{\hat u}|\le\tfrac12$, the worst case is $a/2\le\tfrac12$, which is
\eqref{eq:c2}. Cap against pole gives $b\ip{z}{w}$, which is \eqref{eq:c6}, and
pole against pole is \eqref{eq:c7}.

The remaining four are cap against cap, and we take them by how the two caps
share a line. On one line and one direction the two signs give
$-t+(1-t)=1-2t$, which is at most $\tfrac12$ because \eqref{eq:c3} rearranges to
$t\ge1/(2(1-\gamma))\ge\tfrac12$. On one line and distinct directions the value
is $\pm t+(1-t)\ip{z}{z'}$, worst at $s=s'$, which is \eqref{eq:c4}. On distinct
lines of one $C_i$ it is $\pm t\ip{\hat u}{\hat u'}+(1-t)\ip{z}{z'}$ with
$|\ip{\hat u}{\hat u'}|\le\gamma$; at $z=z'$ that is at most $t\gamma+(1-t)$,
which is \eqref{eq:c3}, and at $z\ne z'$ it is at most
$t\gamma+(1-t)\ip{z}{z'}\le t\gamma+\tfrac12-t<\tfrac12$ by \eqref{eq:c4}, so
this case is implied by the two before it. On lines from distinct $C_i$ and
$C_j$ only $|\ip{\hat u}{\hat u'}|\le\tfrac12$ is available, by
Lemma~\ref{lem:shell}, and $t/2+(1-t)\ip{z}{z'}\le\tfrac12$ is \eqref{eq:c5}.
\end{proof}

\subsection{Pairs versus triples}\label{ss:pairs}

Constraints \eqref{eq:c3} and \eqref{eq:c4} oppose one another. Rearranged,
\eqref{eq:c3} becomes
\begin{equation}\label{eq:tmin}
  t\;\ge\;\frac{1}{2(1-\gamma)},
\end{equation}
so a larger $\gamma$ forces a larger level. Substituting the smallest admissible
$t$ into \eqref{eq:c4} bounds the inner products among directions attached to one
independent set.

\begin{corollary}\label{cor:pairs}
At $t=1/(2(1-\gamma))$ the directions attached to one independent set satisfy
$\ip{z}{z'}\le\gamma/(2\gamma-1)$. Hence
\[
\begin{array}{llll}
  \gamma=\tfrac14 \ (\Leech,\ \Gsz): & t\ge\tfrac23, &
     \ip{z}{z'}\le-\tfrac12, & |Z_i|\le3,\\
  \gamma=\tfrac13 \ (\Pff): & t\ge\tfrac34, &
     \ip{z}{z'}\le-1, & |Z_i|\le2 .
\end{array}
\]
\end{corollary}

\begin{proof}
Substituting $t=1/(2(1-\gamma))$ into $(\tfrac12-t)/(1-t)$ gives
$\gamma/(2\gamma-1)$. If $m$ unit vectors are pairwise at $\ip{z}{z'}\le c<0$
then $0\le|\sum z_i|^2\le m+m(m-1)c$, so $m\le1-1/c$. With $c=-\tfrac12$ this
is $m\le3$ and with $c=-1$ it is $m\le2$.
\end{proof}

Corollary~\ref{cor:pairs} forces the ceiling $\tfrac13$ in \eqref{eq:gamma}: at a
larger threshold the right-hand side of \eqref{eq:c4} falls below $-1$ and no
independent set carries more than one direction. The ceiling is active for $\Gsz$,
where the largest threshold below $\tfrac12$ that Lemma~\ref{lem:ips} leaves
available is $\gamma=\tfrac38$, that is $|\ip{u}{u'}|\le3$; the lemma itself
permits $|\ip{u}{u'}|\le4$, but $\gamma=\tfrac12$ makes \eqref{eq:c3} require
$t\ge1$.
The threshold $\tfrac38$ forces $t\ge\tfrac45$, at which \eqref{eq:c4} requires
$\ip{z}{z'}\le-\tfrac32$, which no two unit vectors satisfy.

Over $\Leech$ and $\Gsz$ an independent set may therefore carry a \emph{zero-sum
triple} of directions: three unit vectors summing to zero and pairwise at
$120^\circ$. Over $\Pff$ it carries at most an antipodal pair. By
\eqref{eq:count} a triple contributes $4$ points per line and a pair $2$,
doubling the gain; no independent set carries more than three directions.

Constraint \eqref{eq:c6} behaves differently at the two levels used below. At
$t=\tfrac34$ we have $b=\tfrac12$, so \eqref{eq:c6} holds for every $z$ and $w$;
$W$ may then be any kissing configuration of $\R^k$, giving $|W|=K(k)$. At
$t=\tfrac23$ we have $b=1/\sqrt3$ and \eqref{eq:c6} requires
$\ip{z}{w}\le\sqrt3/2$; every pole then lies at least $30^\circ$ from every cap
direction. A rotated copy of the direction set is automatically a $60^\circ$
code, so only the $30^\circ$ separation is at issue, and the caps around $Z$ occupy
little of the sphere once $k$ is large: at $k=23$ a single $30^\circ$ cap covers
$2.3\times10^{-8}$ of $S^{22}$, so all $93\,150$ of them together cover at most
$2.1\times10^{-3}$ and a Haar-random rotation is expected to lose at most $198$
directions; the rotation used below loses $76$, leaving $93\,074$. We
therefore take $W$ to be a rotated copy of $Z$ with the points that fail
\eqref{eq:c6} removed. The rotation is made rational by the Cayley transform
$Q=(I-S)(I+S)^{-1}$, which is an isometry of the form whenever $CS$ is
skew-symmetric, so \eqref{eq:c6} is checked in exact integers. It must be an
isometry of $\R^k$ itself and not merely of the coordinates in which $Z$ is
presented, which for some cross-sections are more numerous: $Z$ spans $21$
dimensions of the $24$ that carry $\Lambda_{21}$. Where they differ we take
$S=\sum_t\bigl(z_t(Cz'_t)^{\mathsf T}-z'_t(Cz_t)^{\mathsf T}\bigr)$ over pairs
$z_t,z'_t\in Z$ with $\ip{z_t}{z'_t}=0$, all $2r$ of them pairwise orthogonal.
Then $S$ is integral, $CS$ is skew, and $S$ annihilates the orthogonal complement
of $Z$, so $Q$ fixes that complement pointwise and is an isometry of the span. On
the plane of $z_t$ and $z'_t$, taken in the integral coordinates where
$\ip{z_t}{z_t}=m$, scaling $S$ by $1/\lambda$ makes $Q$ the rotation through
$2\arctan(m/\lambda)$, and its denominator is $\lambda^2+m^2$ for every plane at
once --- two to four digits, against the $2^{71}$ that a basis of the span
would force. Where the direction set is
realised over a ring in which we have not done this we take $W=\emptyset$.

\subsection{Calibration}\label{ss:calib}

We first check \eqref{eq:count} in dimensions where the answer is known. Take
$L=\Leech$, so that $\gamma=\tfrac14$, and take every $C_i$ to be a maximum known
independent set, of $248$ lines. Take $Z$ to be the root system of rank $k$
normalised to the sphere, partitioned into zero-sum triples where its size
permits and into antipodal pairs otherwise; take $|W|=K(k)$ poles. The level
is the smallest admissible one, $t=\tfrac23$, in every column but the first: at
$k=1$ the sphere $S^0$ contains no triple; \S\ref{ss:d25} raises the level to
$t=\tfrac34$ there. Where $t=\tfrac23$ we take the poles to satisfy the
$30^\circ$ separation of \eqref{eq:c6}. Such a set of $K(k)$ poles is exhibited
here in every column. At $k=2$ and $k=4$ the poles are vectors of norm $2m$ in the
same root lattice, where $|a\pm w|^2$ being a lattice norm forces the $30^\circ$
separation and $K(k)$ of them are available. At $k=3$ they are the axis layer of
the dimension-$27$ configuration of \S\ref{ss:d27}, verified over all of its
pairs. At $k=5$, $k=6$ and $k=7$ they are a rotated copy of the root system,
constructed as in \S\ref{ss:big} and checked in exact integers there. So no
column of Table~\ref{tab:calib} assumes a pole set: each is exhibited and
verified. Then
\eqref{eq:count} gives Table~\ref{tab:calib}.

\begin{table}[ht]
\caption{The total \eqref{eq:count} over $\Leech$ with independent sets of $248$
lines, against the published values \cite{cohn-table}. In the third row
$a{\times}b$ means $a$ parts of $b$ directions each, so that $\sum ab=K(k)$; a
part of $3$ is a zero-sum triple and a part of $2$ an antipodal pair. The
columns that disagree are treated in \S\ref{ss:d25} and \S\ref{ss:d27}.}
\label{tab:calib}
\small\setlength{\tabcolsep}{4.5pt}
\begin{tabular}{lrrrrrrr}
\toprule
$k$ & $1$ & $2$ & $3$ & $4$ & $5$ & $6$ & $7$\\
dimension & $25$ & $26$ & $27$ & $28$ & $29$ & $30$ & $31$\\
\midrule
parts of $Z$ & $1{\times}2$ & $2{\times}3$ & $4{\times}3$ & $8{\times}3$ & $12{\times}3{+}2{\times}2$ & $24{\times}3$ & $42{\times}3$\\
\eqref{eq:count} & $197\,058$ & $198\,550$ & $200\,540$ & $204\,520$ & $209\,496$ & $220\,440$ & $238\,350$\\
published & $197\,056$ & $198\,550$ & $200\,044$ & $204\,520$ & $209\,496$ & $220\,440$ & $238\,350$\\
\midrule
& $+2$ & $=$ & $+496$ & $=$ & $=$ & $=$ & $=$\\
\bottomrule
\end{tabular}
\end{table}

The formula reproduces five independently obtained record configurations exactly.
The columns $k=1$ and $k=3$ exceed the published value; \S\ref{ss:d25} and
\S\ref{ss:d27} show that the additional points are available.

\subsection{Dimension 25: the choice of level}\label{ss:d25}

\begin{theorem}\label{thm:25}
$K(25)\ge197\,058$.
\end{theorem}

\begin{proof}
Take $L=\Leech$, $k=1$, one independent set $C$ of $248$ lines, and
$Z=W=\{+1,-1\}\subset S^{0}$. The published configurations use $t=\tfrac23$, at
which $b=1/\sqrt3=0.577\ldots>\tfrac12$; then \eqref{eq:c6} fails for the two
poles, which are discarded.

The level $t=\tfrac23$ is not forced here. Constraint \eqref{eq:tmin} bounds it
from below and \eqref{eq:c4} from above: for two antipodal directions the latter
requires $-1\le(\tfrac12-t)/(1-t)$, that is $t\le\tfrac34$. At $k=1$
the sphere $S^0$ has two points and admits no triple; the whole range
$\tfrac23\le t\le\tfrac34$ is therefore admissible.

At $t=\tfrac34$ constraint \eqref{eq:c3} becomes $\tfrac34\gamma+\tfrac14
\le\tfrac12$, i.e.\ $\gamma\le\tfrac13$; and for $\Leech$ the thresholds
$\gamma=\tfrac14$ and $\gamma=\tfrac13$ define the same condition, since
Lemma~\ref{lem:ips} allows only $\ip{u}{u'}/\mu\in\{0,\pm\tfrac14,\pm\tfrac12\}$
and no admissible value lies strictly between $\tfrac14$ and $\tfrac13$. The
independent set is therefore unchanged at $248$ lines. At $t=\tfrac34$ we have
$b=\tfrac12$, so \eqref{eq:c6} holds and both poles are admissible. Finally
\eqref{eq:count} gives
$196\,560+2\cdot248\cdot(2-1)+2=197\,058$.
\end{proof}

Three constraints meet at $t=\tfrac34$. The bound $t\ge\tfrac23$ keeps the
independent set at $248$ lines, the bound $t\le\tfrac34$ keeps the antipodal pair
admissible, and the poles require $t\ge\tfrac34$. The argument does not repeat for
$k\ge2$: triples exist there and contribute twice as much as pairs, which makes
$t=\tfrac23$ optimal; the poles already attain $K(k)$ at that level.

The gain is $+2$ over any independent set used in its place.

\emph{Superseded, 2026-09-07.} The dimension-$25$ entry of the repository
no longer rests on this level argument. In joint work in progress with H.~Cohn
and B.~Lindow the caps are rebuilt: $1006$ heads $x$ of squared length $3$
(norm-$4$ units) sit in the lens of a minimal vector, each removing exactly its
owner from the equator, and it is a theorem that two compatible heads never
share a removal, so the count is $196\,560+H+2+|E|$ with $H=1006$ and one
further equator point $-\tfrac{2}{\sqrt6}v$ that is not a lattice vector. That
gives $K(25)\ge197\,569$, verified in exact arithmetic by the package
\texttt{dim25-lens-heads}; the account is a separate note. The level argument
above remains the template's own bound and is kept as such.
The construction
uses the $248$-line independent set of \cite{ma-etal2025} unchanged; one of
$249$ lines would give $197\,060$ by the same argument.

\subsection{Dimension 27: partitioning the directions}\label{ss:d27}

\begin{theorem}\label{thm:27}
$K(27)\ge200\,540$.
\end{theorem}

\begin{proof}
Take $L=\Leech$, $k=3$ and $t=\tfrac23$. Take $Z$ to be the $A_3$ root system
normalised to the sphere, the $12$ vertices of a cuboctahedron and a kissing
configuration of $\R^3$, and take $W$ to be a rotated copy of it satisfying the
$30^\circ$ separation of \eqref{eq:c6}, so that $|W|=K(3)=12$. By
Corollary~\ref{cor:pairs} each independent set may carry a zero-sum triple. The
$12$ vertices admit a partition into four zero-sum triples; four independent sets
therefore suffice, and \eqref{eq:count} gives
\[
  196\,560+2\cdot 4\cdot 248\cdot(3-1)+12
  \;=\;196\,560+3\,968+12\;=\;200\,540 . \qedhere
\]
\end{proof}

The choice of direction set is constrained, and not merely convenient. The
$12$-point kissing configuration of $\R^3$ is not unique; two of them are classical,
the cuboctahedron and the icosahedron. Icosahedral vertices meet at
$\ip{z}{z'}\in\{\pm1,\pm1/\sqrt5\}$, so
no two of them that are not antipodal reach the $-\tfrac12$ that
Corollary~\ref{cor:pairs} demands: the icosahedron admits no zero-sum triple at
all, and over it the construction would be reduced to six antipodal pairs and
$199\,548$ points. The cuboctahedron has $24$ pairs at $-\tfrac12$ and admits
exactly two partitions into four triples.

The published configuration \cite{ma-etal2025} groups its twelve directions into
parts of sizes $3,3,2,2,2$. That partition uses five independent sets and deletes
$5\times496=2\,480$ vectors from the equator; four triples use four and delete
$1\,984$, returning $496$ vectors. The construction adds no new spherical code:
the Leech vectors, the twelve directions, the twelve poles and four of the five
independent sets are those of the published configuration; only the
assignment of independent sets to directions differs. A script in \cite{repo}
regenerates the configuration from the dimension-$27$ block of \cite{cohn-data}.

The family admits nothing further in this dimension. By \eqref{eq:count} the gain
is $2\sum_i|C_i|(|Z_i|-1)$ with $|Z_i|\le3$ and $\sum_i|Z_i|\le|Z|=12$, so
$\sum_i(|Z_i|-1)$ is at most $\lfloor 2\cdot12/3\rfloor=8$, which four triples
attain.

\subsection{Dimension 38: large codimension}\label{ss:d38}

For $k\le7$ the number of parts of $Z$ binds in \eqref{eq:count}, and a part is
worth most when it is a triple: the number of triples available,
$\lfloor K(k)/3\rfloor$, is $0,2,4,8,13,24,42$, so at most $42\times248$ lines
carry caps; one therefore takes the independent sets as large as possible. At larger
$k$ the parts of $Z$ outnumber the disjoint maximum independent sets available,
and the binding question becomes how much of the shell admits a partition into
independent sets, a colouring question.

\begin{theorem}\label{thm:38}
$K(38)\ge591\,612$.
\end{theorem}

\begin{proof}[Proof sketch]
Take $L=\Leech$, $k=14$ and $t=\tfrac23$. For $Z$ take the $1\,932$-point kissing
configuration of $\R^{14}$, which admits a partition into $644$ zero-sum triples.
The $98\,280$ minimal lines of $\Leech$ admit a partition into $644$ pairwise
disjoint independent sets, so every line carries caps and the equator is empty:
\eqref{eq:count} gives $196\,560-2\cdot98\,280+2\cdot98\,280\cdot3
=6\cdot98\,280=589\,680$, the ceiling $3K(24)$ of the family.

The pole set is non-empty even though $Z$ already occupies every direction of the
$1\,932$-point code. By \eqref{eq:c6} a pole needs $30^\circ$ of separation from
the directions; the covering radius of that code is $54.7^\circ$. There is a
rotated copy of the code no point of which lies within $30^\circ$ of $Z$, so all
$1\,932$ of them are poles and the total is $589\,680+1\,932=591\,612$, which is
exactly the family's ceiling $3K(24)+K(14)$.

Such a rotation is not found by sampling. A $30^\circ$ cap has relative measure
$1.50\times10^{-5}$ on $S^{13}$, so the $1\,932$ caps cover at most $2.89\%$ of
it; a random rotation accordingly loses $56$ points on average, with standard
deviation $11$, and an independent-points model gives $e^{-56}$ for a clean one.
It is found instead by descending on $\mathrm{SO}(14)$. Scale the $w_i$ to
$|w_i|^2=8$, so that the separation asks for $\ip{Qw_i}{w_j}\le4\sqrt3$; the
objective
$f(Q)=\sum_i\operatorname{softplus}\bigl(\operatorname{smax}_j\ip{Qw_i}{w_j}
-4\sqrt3\bigr)$ is then differentiable, $\operatorname{smax}$ being a smooth
maximum, and the geodesic step
$Q\mapsto\exp(-\eta\,\operatorname{skew}(\nabla f\,Q^{\mathsf T}))Q$ keeps the
iterate on the group. The real solution is
then made rational by the Cayley transform, $(I-S)(I+S)^{-1}$ being rational and
orthogonal for every rational skew-symmetric $S$, so the verification stays
exact: $Q=N/D$ with $N$ integral, $NN^{\mathsf T}=D^2I$ and
$|\ip{Nw_i}{w_j}|\le\operatorname{isqrt}(48D^2)$ for all $1\,932^2$ pairs.
\end{proof}

The partition into independent sets is the substantial step. Let
$b(x,y)=\ip{x}{y}\bmod2$. Since $L$ is even, $b$ descends to an alternating
bilinear form on the $\mathbb F_2$-space $L/2L$. For a totally isotropic subspace
$R\subset L/2L$ of dimension $4$ put
\[
  X_R=\{\,\pm u \;:\; u\in\Lambda,\ b(u,x)=0 \text{ for every } x\in R\,\},
\]
which for the subspaces used here contains $6\,120$ of the $98\,280$ lines. Every
known maximum independent set of $\Gamma_\gamma(\Leech)$ lies inside such a set. Local search restricted to one
$X_R$ returns independent sets of $230$ to $244$ lines, whereas the same search
on the whole graph returns $152$. The explicit partition is in \cite{repo}.

\subsection{Dimensions 49--61 and 73--95}\label{ss:big}

Above dimension $48$ the same construction runs over $\Pff$ and $\Gsz$. Two
things change. First, by Corollary~\ref{cor:pairs}, $\Pff$ admits only antipodal
pairs while $\Gsz$ admits triples; this is the one structural difference between
the two ranges. Second, the construction now needs many disjoint independent
sets; the gain is proportional to their total size.

\begin{proposition}\label{prop:classes}
In the one-point counts of Proposition~\ref{prop:onepoint}, the graph
$\Gamma_\gamma(\Leech)$ has $98\,280$ vertices and is regular of degree
$n[2]=4\,600$. The graph $\Gamma_\gamma(\Pff)$ has $26\,208\,000$ vertices and is
regular of degree $n[3]=36\,848$. The graph $\Gamma_\gamma(\Gsz)$ has
$3\,109\,087\,800$ vertices and is regular of degree $n[3]+n[4]=14\,086\,968$. By
the Caro--Wei bound they contain independent sets of size at least $22$, $712$
and $221$.
\end{proposition}

Search on explicit coordinates finds independent sets of $248$ lines in the
first of these graphs, $7\,069$ in the second and $2\,118$ in the third. The
ratio of the size found to the Caro--Wei bound is $11.3$, $9.93$ and $9.58$.

Two ways of producing a disjoint family are used. The first takes automorphisms:
if $g\in\mathrm{Aut}(L)$ and $C$ is independent then so is $Cg$, and two random
images of a set of size $c$ overlap in about $c^2/\#\{\text{lines}\}$ lines,
which we delete from the later image. This gives $1\,400$ pairwise disjoint
independent sets of $\Pff$, covering $8\,023\,414$ lines or $30.6\%$ of the
minimal lines, and $32\,000$ of $\Gsz$, covering $47\,150\,230$ lines. The second
runs the search itself on a pool of minimal lines, removing each set from the
pool as it is found; over $\Gsz$ it gives $32\,000$ sets covering
$65\,570\,602$ lines. The two have opposite shapes. The images of one good set
begin at $2\,106$ lines and decay to $867$, while the sets found afresh lie
between $1\,932$ and $2\,118$. A dimension that reads few sets is better served
by the first family and one that reads many by the second, the crossover falling
at dimension $86$. The two are never mixed, disjointness across families not
being claimed.

\begin{theorem}\label{thm:big}
For $49\le n\le61$ and $73\le n\le95$ the lower bounds of
Table~\ref{tab:all} hold.
\end{theorem}

\begin{proof}[Proof sketch]
Apply \eqref{eq:count} with $L=\Pff$, $t=\tfrac34$ and $|Z_i|=2$ for
$49\le n\le61$. For $73\le n\le95$ apply it with $L=\Gsz$, taking $t=\tfrac23$
and $|Z_i|=3$ where a partition of $Z$ into zero-sum triples is available, and
$t=\tfrac34$ with $|Z_i|=2$ otherwise. The direction set $Z$ is a root system for
$k\le8$ and a cross-section of $\Leech$ for larger $k$; if its partition has $r$
parts, take the $r$ largest available independent sets. At $t=\tfrac34$ the poles
are unconstrained and contribute $K(k)$; at $t=\tfrac23$ we take for $W$ a rotated
copy of $Z$ restricted to the points satisfying the $30^\circ$ separation of
\eqref{eq:c6}, which is between $99\%$ and $100\%$ of it;
it is all of $Z$ at $k=3$, $5$, $6$, $7$, $8$, $9$, $10$, $12$, $13$, $14$, $15$ and
$16$,
and reaches $K(k)$ at $k=3$, $5$, $6$, $7$, $8$, $9$, $10$, $14$ and $16$.

The choice of $Z$ is not forced. Conditions \eqref{eq:c4} and \eqref{eq:c5} ask of it
only that it be a $60^\circ$ code admitting a partition into parts internally at cosine
$\le-\tfrac12$, so a larger such code with a perfect partition is strictly better, and at
$k=12$ and $13$ the cross-section taken was not the largest available. Let $\pi$ be an
automorphism of the Golay code of cycle shape $1^63^6$, acting on $\Leech$ by permuting
coordinates. Its fixed space is $12$-dimensional --- six fixed points and six $3$-cycles ---
and on the orthogonal complement, also $12$-dimensional, $\pi$ has no eigenvalue $1$, so
$I+\pi+\pi^2=0$ there and every orbit $\{v,\pi v,\pi^2v\}$ of a vector of that complement
sums to zero. The $756$ minimal vectors of $\Leech$ lying in the complement have
determinant $3^6$ at minimal norm $4$, which is the Coxeter--Todd lattice $K_{12}$, and
their partition into $252$ zero-sum triples is the orbit set of $\pi$: a group acting, not
a packing. Adjoining the one further direction carrying the most minimal vectors --- $162$
of them --- gives a $13$-dimensional section of $918$, against $\Lambda_{12}$'s $648$ and
$\Lambda_{13}$'s $906$. Thirteen is odd, so no order-$3$ isometry of $\R^{13}$ acts freely
and those $306$ triples are packed rather than given.

The same identity settles four partitions that a packing search had left short. An isometry
$S$ with $I+S+S^2=0$ has eigenvalues $\omega$ and $\bar\omega$ only, so it acts freely and
all its orbits are zero-sum triples; an $S$-invariant $Z$ is therefore partitioned perfectly
by them, and those eigenvalues come in conjugate pairs, so this can happen only in even
dimension. At $k=16$, $18$, $20$ and $22$ such an $S$ exists and the partitions improve
from $1\,434$ triples and $7$ pairs to $1\,440$ triples, from $2\,457$ and $11$ to $2\,466$,
from $5\,782$ and $23$ to $5\,800$, and from $16\,594$ and $49$ to $16\,632$. At
$k=20$ the isometry condition alone leaves the first level of the search $664$ ways
wide and a second invariant is needed: $S$ carries zero-sum triples to zero-sum triples,
so it preserves the number of them through a vector, and filtering the candidate images by
that degree narrows the level to eight. At $k=22$ that degree is constant and the filter
buys nothing; the isometry condition settles it there on its own, in twelve complete
assignments. The six dimensions this moves are $84$, $85$, $88$, $90$, $92$ and $94$, by
$549\,638$ spheres in total.

The odd $k$ carry no such $S$, and their partitions had been left short by a greedy that
never revisits a placed triple: $9$, $13$, $25$ and $60$ triples at $k=17$, $19$, $21$ and
$23$. What closes them is that the hypergraph of zero-sum triples is \emph{linear} --- two
vectors lie in at most one triple, since the third is forced to be $-(a+b)$ --- which
determines the move set of a local search completely. If a triangle $\{u,v,w\}$ has $u$
uncovered and $v$, $w$ in a common part, that part's third vector is $-(v+w)=u$, which is
uncovered; so it does not happen. Evicting one part therefore always leaves the number of
uncovered vectors unchanged, and every neutral or improving move is a pair of uncovered
vectors at $120^\circ$. Enumerating those pairs instead of sampling triangles takes the
partitions from $1\,773$ triples and $11$ pairs to $1\,782$, from $3\,543$ and $16$ to
$3\,556$, from $9\,215$ and $34$ to $9\,240$, and from $30\,990$ and $82$ to $31\,050$,
perfect in every case, and moves dimensions $89$, $91$, $93$ and $95$ by $286\,996$
spheres.

The economy that suggests itself is not available. $Z=-Z$ and a triple's signs are fixed up to
a global one, so its triples come in antipodal pairs and the search appears to halve --- and
an $S$-invariant partition is antipodally symmetric for the same reason, the orbit of $-v$
being the negative of the orbit of $v$. Summing $[a]+[b]+[c]=0$ in $\Lambda/2\Lambda$ over
such a partition gives $\sum[v]=0$ over the $|Z|/2$ minimal lines, and for $\Lambda_k$ that
fails at $k=11$, $13$, $15$, $17$, $19$ and $23$. On the vectors the same sum is doubled and
says nothing, which is why the search runs there; a search on the lines stalls one triangle
short at $k=17$, at $890$ of $891$, for arithmetic reasons rather than for want of looking.

A copy of $Z$ cannot exceed $|Z|$, and at $k=19$, $20$ and $21$ the direction set is the
minimal-vector set of $\Lambda_k$, which is smaller than $K(k)$: $10\,668$, $17\,400$ and
$27\,720$ against $11\,948$, $19\,448$ and $29\,768$. Those three layers were short by
choice of ansatz rather than by geometry, since \eqref{eq:c6} and \eqref{eq:c7} ask of $W$
only that it be a $60^\circ$ code lying $30^\circ$ from $Z$ --- not that it be a copy of
anything. Taking for $W$ a rotated copy of the record configuration of $\R^k$ instead gives
$11\,934$, $19\,415$ and $29\,712$ poles. At $k=19$ and $20$ that configuration is available
in the coordinates the direction set is already written in, $\Lambda_k$ there being
$4\binom{n}{2}$ vectors $(\pm2,\pm2,0,\dots,0)$ together with the $128$ even-sign patterns
on each octad of a shortened Golay code; taking the octad signs odd and adjoining
$((-1)^{c_1},\dots,(-1)^{c_n})$ for $c$ in the dual of the octad code is Cohn--Li's
construction, and it attains $K(k)$.

At $k=21$ the direction set is written in $24$ coordinates spanning $21$, so a configuration
built in $\R^{21}$ must be carried into the axis space by a rational similarity: $k$
orthogonal vectors of a common norm $\mu$ in the span. Whether one exists is decided by the
discriminant of that space, since an orthogonal basis of norms $\mu$ has determinant $\mu^k$
modulo squares. At $k=18$ and $k=22$ the discriminant is $3$ and $k$ is even, so none exists
at all and no configuration built in the standard $\R^k$ can be placed there rationally; at
$k=21$ it is $1$, so $\mu$ must be a square, which is why a frame drawn from the norm-$32$
direction vectors is always one short and one drawn from norm-$4$ vectors of the span is not.

That is a statement about the standard metric and not about the configuration. Three further
direction sets are cross-sections written in more coordinates than they span, at $k=15$,
$17$ and $18$, and there the record configuration is not rebuilt but read: Cohn's
coordinate data set gives it as exact integers together with an integer metric
$\operatorname{diag}(c)$, which for $k=15$, $17$, $18$ is
$\operatorname{diag}(1^{14},2)$, $\operatorname{diag}(1^{16},2)$ and
$\operatorname{diag}(1^{16},2,6)$. A frame then has to satisfy
$M^{\mathsf T}M=\mu\operatorname{diag}(c)$ and the discriminant condition becomes
$\mu^k\det\operatorname{diag}(c)\equiv\operatorname{disc}$. At $k=18$ that is
$12\equiv3$, the discriminant of the axis space, and a frame does exist: sixteen columns
$v_1\pm v_2,\dots,v_{15}\pm v_{16}$ of norm $64$ built from mutually orthogonal direction
vectors, $2v_{17}$ of norm $128$, and for the last one the orthogonal complement of those
seventeen inside the span, which is a line and whose norm class is therefore forced --- it
comes out at $6$ modulo squares, which is what the frame needs. The rotation is taken to be
a product of rational plane rotations of $\operatorname{diag}(c)$, each of denominator $D$
with $a^2+(c_j/c_i)b^2=D^2$, since a full-rank Cayley transform on a space of this rank has
denominator about $d^k$ and cannot be rounded finely enough to spend the margin. The layers
go from $2\,340$, $5\,338$ and $7\,374$ to $2\,562$, $5\,704$ and $7\,624$ poles, against $K(k)$ = $2\,564$, $5\,730$ and $7\,654$.

Dimension $83$ is no longer an exception either. Its configuration is realised over
$\Z[\sqrt2]$, but only the configuration is: a rational $Q=N/D$ with $Q^{\mathsf T}CQ=C$
carries the directions to poles in the same ring, and \eqref{eq:c6} becomes
$(p+q\sqrt2)^2\le972D^2$ for integers $p,q$, decided exactly. All $604 = K(11)$ of them
clear.
\end{proof}

Every dimension from $81$ to $95$ now carries its own construction: dimension $81$
exceeds dimension $80$ by $185\,042$, so no entry of Table~\ref{tab:all} in this
range rests on monotonicity of $K$. The construction also runs in dimensions $62$
and $63$, where it reaches $64\,217\,822$ and $67\,365\,752$ and is beaten;
\S\ref{ss:d6263} treats those two.

The layered construction stops winning at $96$. Throughout $n\le95$ the
Edel--Rains--Sloane total of \S\ref{sec:ers} stays below $\Gsz$. Its largest term
is $A(n_0,\lceil n_0/4\rceil)$, for which the code tables give at most $2^{30}$
up to $n_0=92$, or $17\%$ of $6\,218\,175\,600$. At $n_0=93,94,95$ they give
$2^{31}$, $2^{31}$ and $2^{32}$, and the whole chain evaluated at those three
lengths gives $2.68\times10^9$, $2.68\times10^9$ and $4.83\times10^9$. At $n_0=96$ the
code $[96,33,24]$ makes the leading term $A(96,24)\ge2^{33}$, at which the chain
$(96,24,6,1)$ totals $12\,886\,999\,232$, against the $6\,480\,558\,568$ the
layered construction reaches there. Dimension $96$ is therefore
Edel--Rains--Sloane's rather than $\Gsz$'s, and \S\ref{ss:d96} takes it on that
footing --- as \S\ref{ss:d6263} does dimensions $62$ and $63$.

Every input to those evaluations is a construction, so each is a lower bound for
the Edel--Rains--Sloane count. A lower bound above $\Gsz$ settles which of the
two constructions wins there; a lower bound below it settles nothing by itself.
\S\ref{ss:exposure} measures the margin in the dimensions above $48$.

\emph{And the table stops there.} $K$ is superadditive: two $60^\circ$ codes placed on
orthogonal subspaces meet at inner product $0$, so $K(96+m)\ge K(96)+K(m)$, and every
dimension above $96$ already follows from the row for $96$ together with a published
low-dimensional record. A row for dimension $100$ reading $12\,886\,999\,232+24$ would be
arithmetic on the row above it. What that leaves open is whether anything overtakes the
implied floor above $96$, as Edel--Rains--Sloane overtakes $\Gsz$ \emph{at} $96$. The layered
construction does not: it ends at $k=24$. The chain does, but by a margin that decides
nothing --- evaluated through $n=128$ from the same tables used here, it exceeds the implied
floor by at most $0.016\%$, at $n=114$. So $96$ is where the table stops saying anything, which is why it
stops there.

\section{The Edel--Rains--Sloane chain}\label{sec:ers}

Four improvements are code-theoretic. The Edel--Rains--Sloane construction
\cite{edel-rains-sloane1998} builds a $60^\circ$ code from a
chain $n\ge n_0\ge4n_1\ge16n_2\ge\cdots$ of support sizes, giving
\[
  K(n)\;\ge\;\sum_\nu A(n,n_\nu,n_\nu)\,A\bigl(n_\nu,\lceil n_\nu/4\rceil\bigr).
\]
The construction is defined for every $n$, but its authors evaluate it only at
$n=32,36,40,44,64,80,128$. Elsewhere its value is a question about the current
code tables rather than about the construction.

\subsection{Dimension 39}\label{ss:d39}

In dimensions $38$ and $39$ the chain that maximises the total has three
levels, $(n,8,2)$, and reads
\[
  K(n)\;\ge\;A\bigl(n,\lceil n/4\rceil\bigr)+128\,A(n,8,8)+4\binom{n}{2}.
\]

\begin{theorem}\label{thm:39}
$K(39)\ge756\,116$.
\end{theorem}

\begin{proof}
Echols \cite{echols2026} improves $A(39,8,8)$ from $3\,323$ to $3\,324$;
Brouwer's table \cite{brouwer} carries the new value. With
$A(39,10)\ge327\,680$ \cite{zinoviev-litsyn1984},
\[
  327\,680+128\cdot3\,324+4\cdot741 \;=\; 327\,680+425\,472+2\,964 \;=\; 756\,116 . \qedhere
\]
\end{proof}

The gain is $128\,(3324-3323)=128$. Echols \cite{echols2026} computes kissing
numbers from the fixed level-one length $n_0=32$, which in dimension $39$ falls
below the record; that paper therefore reports no value there. The chain
optimises at $n_0=n$ in dimensions $38$, $39$, $40$, $43$ and $44$; among these, $38$ and $39$
are the two whose totals exceed the published entry; \S\ref{ss:d38} supersedes
dimension $38$. That entry is $755\,988$, which
\cite{cohn-table} attributes to the Edel--Rains--Sloane construction; of the three
dimensions named in \cite{sun-wang2026} it is dimension $43$, not dimension $39$, whose
entry that table credits to the antipode construction. That paper also reconstructs the
twelve cross-section packings of $1982$ in dimensions $38$--$47$ and computes their
kissing numbers, reporting that they exceed the best previously tabulated lower bounds
in dimensions $42$--$47$. Dimensions $38$ and $39$ are the only two this paper claims
inside that span of $38$--$47$, and neither lies in $42$--$47$.

The verification in \cite{repo} checks the constant-weight code over all
$\binom{3324}{2}$ pairs and the ten geometric conditions of the chain as exact
rational inequalities, three of them tight. The value $A(39,10)\ge327\,680$ is
taken from \cite{zinoviev-litsyn1984} and is not verified here; the existing
record $755\,988$ rests on the same value.

\subsection{Dimensions 62 and 63}\label{ss:d6263}

The first term of the chain needs no chain. A binary code of length $n_0$, dimension
$m$ and minimum distance $d$ sends each codeword $c$ to the unit vector
\[
  x(c)\;=\; n_0^{-1/2}\bigl((-1)^{c_1},\dots,(-1)^{c_{n_0}}\bigr),
\]
and two codewords at Hamming distance $j$ satisfy
$\ip{x(c)}{x(c')}=(n_0-2j)/n_0$. Every pair is therefore at $60^\circ$ or more
exactly when $d\ge\lceil n_0/4\rceil$; the $2^m$ sign vectors are then a
kissing configuration of $\R^{n_0}$, hence of $\R^n$ for every $n\ge n_0$:
\[
  K(n)\;\ge\; \max_{n_0\le n} A\bigl(n_0,\lceil n_0/4\rceil\bigr).
\]

Since $\lceil 62/4\rceil=\lceil 63/4\rceil=16$, the codes $[62,26,16]$ and
$[63,27,16]$ of \cite{codetables} qualify, giving $2^{26}$ and $2^{27}$ sign
vectors. Both are shortenings of one $[64,28,16]$; shortening an $[n,m,d]$ code
gives $[n-1,m-1,\ge d]$. That alone already beats the values of \S\ref{ss:tables}.

The rest of the chain is not negligible: it carries Edel--Rains--Sloane's own
dimension-$64$ total from $2^{28}=268\,435\,456$ to $331\,737\,984$. Since
$4\cdot16=64>63$, the largest level-one weight the chain allows in either dimension
is $n_1=15$, and $15\cdot4=60\le62$, so the chain is $(n,15,2)$ and
\[
  K(n)\;\ge\;A\bigl(n,\lceil n/4\rceil\bigr)+A(15,4)\,A(n,15,15)+4\tbinom{n}{2}.
\]

\begin{theorem}\label{thm:6263}
$K(62)\ge71\,310\,732$ and $K(63)\ge138\,419\,844$.
\end{theorem}

\begin{proof}
Level zero gives $2^{26}$ and $2^{27}$ as above. For level one, index the cells of a
$15\times4$ grid by $60$ of the coordinates and let a support take one cell in each
row; two supports then meet in exactly $15-d_H$ of their choice words, so a
quaternary code of length $15$ and minimum distance $8$ yields supports meeting in
at most $7=\lfloor15/2\rfloor$. The cyclic $[15,6,8]_4$ with defining set
$\{0,1,2,3,4,5,8,10,12\}$, a union of cyclotomic cosets modulo $15$ under
$x\mapsto4x$, gives $A(n,15,15)\ge4^6=4\,096$, and $A(15,4)=1\,024$ from the
extended Hamming $[16,11,4]$ shortened once. Level two takes every pair of
coordinates with all four sign patterns. Hence
\begin{align*}
  67\,108\,864+4\,194\,304+7\,564 &= 71\,310\,732,\\
  134\,217\,728+4\,194\,304+7\,812 &= 138\,419\,844 . \qedhere
\end{align*}
\end{proof}

Everything above level zero is built rather than cited: \cite{repo} enumerates all
$4^6$ codewords of the quaternary code to establish its minimum distance, checks all
$\binom{4096}{2}$ support overlaps directly, enumerates all $2^{10}$ codewords of the
$[15,10,4]$, and runs the whole chain end to end in coordinates at the two dimensions
where it happens to reproduce the exact kissing number --- $n=8$, chain $(8,2)$,
giving the $240$ points of the $E_8$ root system, and $n=16$, chain $(16,4,1)$, giving
$4\,320$, the Barnes--Wall value and the record --- with every pair of both
configurations checked. Only the two level-zero code parameters are taken from
\cite{codetables}.

Against the values of \S\ref{ss:tables} these are factors $1.36$ and $2.64$. Level
zero alone already exceeds the \emph{ceiling} of the construction of
\S\ref{ss:big} in these two dimensions, so no family of independent sets can
recover them there. Over $\Pff$ that construction gives
$52\,416\,000+2\sum_i|C_i|+K(k)$ by \eqref{eq:count} with $|Z_i|=2$. The number
of independent sets is at most the number $\lfloor K(k)/2\rfloor$ of antipodal
pairs available among the directions, and each $|C_i|$ is at most the largest
independent set known, of $7\,069$ lines. At $k=14$ and $k=15$ that gives
$66\,075\,240$ and $70\,543\,480$, short of $2^{26}$ and $2^{27}$ by
$1\,033\,624$ and $63\,674\,248$. Closing those gaps would need independent sets
of $7\,605$ and $31\,903$ lines against the largest known $7\,069$, or more parts
than the record in dimension $14$ or $15$ allows. The chain is not a refinement of
\S\ref{ss:big} in these dimensions but a different and better construction.

\subsection{Dimension 96}\label{ss:d96}

At $n=96$ the chain closes up: $96\ge96\ge4\cdot24\ge16\cdot6\ge64\cdot1$, so
$(96,24,6,1)$ is admissible, and every inequality but the last is an equality.
Level zero is again a single sign code, and here it carries two thirds of the
total: $\lceil96/4\rceil=24$, and the code $[96,33,24]$ of \cite{codetables}
gives $A(96,24)\ge2^{33}$.

The two middle levels apply the grid map of \S\ref{ss:d6263} at two different
factorisations of $96$. Writing $96=24\cdot4$ and indexing the cells of a
$24\times4$ grid by the coordinates, a support taking one cell in each row has
weight $24$, and two supports meet in exactly $24-d_H$ of their choice words; a
quaternary code of length $24$ and minimum distance $12$ therefore gives supports
meeting in at most $12=\lfloor24/2\rfloor$, and $[24,9,12]_4$ of \cite{codetables} gives
$A(96,24,24)\ge4^9$. Writing instead $96=6\cdot16$, the same argument over
$\mathbb F_{16}$ with the Reed--Solomon code $[6,4,3]_{16}$ gives
$A(96,6,6)\ge16^4$.

\begin{theorem}\label{thm:96}
$K(96)\ge12\,886\,999\,232$.
\end{theorem}

\begin{proof}
The four levels contribute
\[
  A(96,24)+A(96,24,24)\,A(24,6)+A(96,6,6)\,A(6,2)+A(96,1,1)\,A(1,1),
\]
in which $A(6,2)=32$, $A(1,1)=2$ and $A(96,1,1)=96$ are exact, while \cite{brouwer}
brackets $A(24,6)$ between $16\,384$ and $24\,106$; a larger value would only raise
the total, so the smaller end is substituted. With that and the three lower bounds
above,
\[
  8\,589\,934\,592+262\,144\cdot16\,384+65\,536\cdot32+96\cdot2
  \;=\;12\,886\,999\,232 . \qedhere
\]
\end{proof}

The value of \S\ref{ss:tables} here is the direct sum $\Gsz\oplus\Leech$, at
$6\,218\,175\,600+196\,560$, so the factor is $2.07$. As in \S\ref{ss:d6263} the
mathematics is Edel--Rains--Sloane's, and as there we find no record of their
construction having been evaluated at this length, which is not among their seven
worked examples; their own first
remark licenses substituting any available lower bounds for the code sizes, which
is all that is done above. The total is a function of those tables: with the three
levels below weight $96$ as constructed here --- the weight-$24$, weight-$6$ and
weight-$1$ terms, together $4\,297\,064\,640$ --- the chain passes the direct sum
as soon as $A(96,24)\ge2^{31}$, and $[96,33,24]$ gives $2^{33}$. Of the three code inputs only the smallest is
built here: \cite{repo} constructs $[6,4,3]_{16}$ and establishes its minimum distance
by evaluating all $16^4$ polynomials. The sign code $[96,33,24]$, which carries two
thirds of the total, and the quaternary code $[24,9,12]_4$ are taken from \cite{codetables}.

\section{The results}\label{sec:results}

\begin{longtable}{rrrl}
\caption{All forty-seven improvements. ``Previously'' is \cite{cohn-table} where
that table has an entry and otherwise the direct sum or inherited value of
\S\ref{ss:tables}. Dimensions $44$ to $47$ are absent:
\cite{sun-wang2026} gives larger values in dimensions $44$ and $45$;
dimensions $46$ and $47$ are already in the literature (\S\ref{ss:new}).}\label{tab:all}\\
\toprule
$n$ & previously & this work & construction\\
\midrule
\endfirsthead
\toprule
$n$ & previously & this work & construction\\
\midrule
\endhead
\bottomrule
\endfoot
$25$ & $197\,056$ & $197\,569$ & layers over $\Leech$, lens heads (\S\ref{ss:d25})\\
$27$ & $200\,044$ & $200\,540$ & layers over $\Leech$, four triples\\
$38$ & $566\,652$ & $591\,612$ & layers over $\Leech$, $k=14$\\
$39$ & $755\,988$ & $756\,116$ & constant-weight codes\\
$49$ & $52\,416\,002$ & $52\,430\,140$ & layers over $\Pff$\\
$50$ & $52\,416\,006$ & $52\,458\,418$ & layers over $\Pff$\\
$51$ & $52\,416\,012$ & $52\,500\,816$ & layers over $\Pff$\\
$52$ & $52\,416\,024$ & $52\,585\,516$ & layers over $\Pff$\\
$53$ & $52\,416\,040$ & $52\,698\,222$ & layers over $\Pff$\\
$54$ & $52\,416\,072$ & $52\,922\,906$ & layers over $\Pff$\\
$55$ & $52\,416\,126$ & $53\,299\,730$ & layers over $\Pff$\\
$56$ & $52\,416\,240$ & $54\,085\,392$ & layers over $\Pff$\\
$57$ & $52\,416\,306$ & $54\,534\,176$ & layers over $\Pff$\\
$58$ & $52\,416\,510$ & $55\,893\,422$ & layers over $\Pff$\\
$59$ & $52\,416\,604$ & $56\,505\,668$ & layers over $\Pff$\\
$60$ & $52\,416\,841$ & $57\,477\,123$ & layers over $\Pff$\\
$61$ & $52\,417\,154$ & $59\,898\,694$ & layers over $\Pff$\\
$62$ & $52\,417\,932$ & $71\,310\,732$ & \textsc{ers} chain $(62,15,2)$\\
$63$ & $52\,418\,564$ & $138\,419\,844$ & \textsc{ers} chain $(63,15,2)$\\
$68$ & $331\,737\,984$ & $361\,275\,480$ & cross-section of $\Gsz$, $k=4$\\
$69$ & $331\,737\,984$ & $627\,822\,180$ & cross-section of $\Gsz$, $k=3$\\
$70$ & $331\,737\,984$ & $1\,249\,778\,250$ & cross-section of $\Gsz$, $k=2$\\
$71$ & $331\,737\,984$ & $2\,603\,658\,750$ & cross-section of $\Gsz$, $k=1$\\
$73$ & $6\,218\,175\,602$ & $6\,218\,179\,814$ & layers over $\Gsz$\\
$74$ & $6\,218\,175\,606$ & $6\,218\,192\,454$ & layers over $\Gsz$\\
$75$ & $6\,218\,175\,612$ & $6\,218\,209\,308$ & layers over $\Gsz$\\
$76$ & $6\,218\,175\,624$ & $6\,218\,243\,016$ & layers over $\Gsz$\\
$77$ & $6\,218\,175\,640$ & $6\,218\,285\,152$ & layers over $\Gsz$\\
$78$ & $6\,218\,175\,672$ & $6\,218\,377\,840$ & layers over $\Gsz$\\
$79$ & $6\,218\,175\,726$ & $6\,218\,529\,454$ & layers over $\Gsz$\\
$80$ & $6\,218\,175\,840$ & $6\,218\,849\,228$ & layers over $\Gsz$\\
$81$ & $6\,218\,175\,906$ & $6\,219\,034\,270$ & layers over $\Gsz$\\
$82$ & $6\,218\,176\,110$ & $6\,219\,605\,578$ & layers over $\Gsz$\\
$83$ & $6\,218\,176\,204$ & $6\,219\,621\,690$ & layers over $\Gsz$\\
$84$ & $6\,218\,176\,441$ & $6\,220\,293\,304$ & layers over $\Gsz$\\
$85$ & $6\,218\,176\,754$ & $6\,220\,745\,414$ & layers over $\Gsz$\\
$86$ & $6\,218\,177\,532$ & $6\,223\,567\,844$ & layers over $\Gsz$\\
$87$ & $6\,218\,178\,164$ & $6\,224\,703\,598$ & layers over $\Gsz$\\
$88$ & $6\,218\,179\,920$ & $6\,230\,206\,492$ & layers over $\Gsz$\\
$89$ & $6\,218\,181\,330$ & $6\,233\,054\,848$ & layers over $\Gsz$\\
$90$ & $6\,218\,183\,254$ & $6\,238\,744\,788$ & layers over $\Gsz$\\
$91$ & $6\,218\,187\,548$ & $6\,247\,801\,034$ & layers over $\Gsz$\\
$92$ & $6\,218\,195\,048$ & $6\,266\,405\,667$ & layers over $\Gsz$\\
$93$ & $6\,218\,205\,368$ & $6\,294\,850\,360$ & layers over $\Gsz$\\
$94$ & $6\,218\,225\,496$ & $6\,355\,739\,600$ & layers over $\Gsz$\\
$95$ & $6\,218\,268\,750$ & $6\,473\,065\,046$ & layers over $\Gsz$\\
$96$ & $6\,218\,372\,160$ & $12\,886\,999\,232$ & \textsc{ers} chain $(96,24,6,1)$\\
\end{longtable}

\section{Remarks}\label{sec:remarks}

\subsection{Obstructions}\label{ss:obstructions}

Several constructions reach their exact ceiling. Details and certificates are in
\cite{repo}.

\emph{The Cohn--Li construction is exhausted in dimensions $18$, $20$ and $21$.}
It adjoins to a base configuration a set of vectors indexed by words orthogonal
to a constant-weight code, two such vectors being compatible exactly when their
Hamming distance is at least $n/4$ \cite{cohn-li2024}. The words in question form
a punctured Golay code of $4\,096$ elements. The number adjoined is therefore the
independence number of a Cayley graph on $\mathbb F_2^{12}$, whose connection set
$\Sigma$ consists of the non-zero words of weight less than $n/4$. For an
abelian Cayley graph the Delsarte bound equals the Lov\'asz theta function. In
dimensions $20$ and $21$ both equal $2\,048$, which \cite{cohn-li2024} attains.
In dimension $20$ the set $\Sigma$ has five elements and they are linearly
independent, so each of the $128$ components of the graph is the $5$-cube,
bipartite with independence number $16$, whence $128\times16=2\,048$. In dimension
$18$ the number is $768$ and is attained, but the construction then totals
$6\,500+768=7\,268$, short of the $7\,654$ that \cite{cohn-li2024} prove in that
dimension by a configuration of a different shape.

\emph{Dimension $19$ is closed at the value of Ho.} Here the graph has four
components. Each is the Cayley graph on $\mathbb F_2^{10}$ with a connection set
$\Sigma$ of $21$ elements; the construction produces $10\,668+4\alpha$
points, where $\alpha$ is the independence number of one component. Exactly one
element $t$ of $\Sigma$ has odd weight. Writing $H$ for the even-weight
hyperplane, the two halves $H$ and $H+t$ induce the same
$512$-vertex graph $G_0$, and the edges between them form the perfect matching
$x\mapsto x+t$; hence $\alpha\le2\alpha(G_0)$. Constraint programming, with the
inequality that an independent set meets a $5$-cycle of $G_0$ in at most two
vertices, proves $\alpha(G_0)=160$ optimal. Hence $\alpha\le320$, which
\cite{ho2026} attains; $11\,948$ is the exact ceiling of the construction.

For a kissing configuration $X\subset S^{d-1}$ put
\[
  m(X)=\min_{z\in S^{d-1}}\ \max_{x\in X}\ \ip{z}{x}.
\]
A further point may be added to $X$ exactly when $m(X)\le\tfrac12$, so $X$ is
maximal exactly when $m(X)>\tfrac12$; we call $2m(X)-1$ the \emph{margin} of
$X$. The values of $m$ below are obtained by numerical maximisation of $|z|$ over
the polytope $\{z:\ip{z}{x}\le\tfrac12 \text{ for all } x\in X\}$, whose maximum
is $1/(2m(X))$; where an exact algebraic value is quoted, it is confirmed in
rational arithmetic.

\emph{$\Lambda_{21}$, $\Lambda_{22}$ and $\Lambda_{23}$ are maximal spherical
codes,} with $m=\sqrt{8/29}$, $\sqrt{3/11}$ and $\sqrt{4/15}$, so no point may be
added to any of them; the level argument of \S\ref{ss:d25} has no analogue
there. The remaining cells of \S\ref{ss:onepoint} add nothing either. Fix
$u\in\Lambda$ and project the shell of $\Leech$ onto $u^\perp$. The cell $n[0]$
already lies in $u^\perp$ and is the minimal shell of $\Lambda_{23}$. Two
projected vectors taken from opposite signs of the cell $n[1]$, or of the cell
$n[2]$, can have inner product above $\tfrac12$; call such a pair a conflict. At
most one sign of each cell is therefore available. On the union of $n[0]$ with
one sign of each of the other two, conflicts occur only between $n[0]$ and the
rest: each vector of $n[0]$ conflicts with $1\,024$ of the $47\,104$ and with
$44$ of the $4\,600$, while these conflict with $2\,025$ and $891$ vectors of
$n[0]$ respectively. Weighting each edge by
$1/2025$ or by $1/891$ accordingly gives a fractional matching under which every
vector of $n[0]$ carries load $1024/2025+44/891<1$. The bipartite matching
polytope is integral, so a matching saturates the $51\,704$ vectors outside
$n[0]$. By K\"onig's theorem the largest kissing configuration obtainable this
way has exactly $93\,150$ points.

\emph{Every published record in dimensions $9$--$19$ is maximal,} its largest cosine
$m$ exceeding $\tfrac12$ by between $4.4\%$ (dimensions $17$ and $19$) and $34.2\%$
(dimension $10$). Dimensions $14$ and $15$ have the same value $m=\sqrt{2/7}$.

\emph{The antipode construction is bounded by a spherical cap.} Fix linearly
independent $u_1,\dots,u_k\in\Lambda$, let $\pi$ denote the orthogonal projection
of $\R^n$ onto their span, and for $p\in\pi(L)$ put
$f(p)=\#\{w\in\Lambda:\pi(w)=p\}$, so that $f(0)=|S|$ in the notation of
Lemma~\ref{lem:cross}. The construction of \cite{sun-wang2026} chooses a finite
$T\subset\pi(L)$ whose largest squared pairwise distance $\delta$ is less than
$\mu$ and returns the maximum over $p\in T$ of $\sum f(q-p)$, the sum running
over the $q\in T$ with $|q-p|^2=\delta$. Over all such $T$ this maximum equals
the largest weight of a clique, taken over all $\delta\in(0,\mu)$, in the graph
on $\{v\in\pi(L):|v|^2=\delta,\ f(v)>0\}$ in which $v$ carries weight $f(v)$ and
$v,v'$ are adjacent when $\ip{v}{v'}\ge\delta/2$. Adjacency
means an angle of at most $60^\circ$, so a clique lies in a spherical cap of that
angular radius about any of its members, and the construction cannot exceed the
total weight of the heaviest such cap. Two members lie within $30^\circ$ of their
bisector, but the radius does not stay there: three members pairwise at exactly
$60^\circ$ need $\arccos\sqrt{2/3}=35.26^\circ$. The cap bound improves on $f(0)$
alone only when $f(0)$ is already within about $10\%$ of the value it has to
exceed, as in dimensions $44$ through $47$.

\subsection{The margin above dimension 48}\label{ss:exposure}

Among the dimensions at issue the tables have a single entry, at $n=80$, where it
agrees with the direct sum of \S\ref{ss:tables}. Everywhere else above dimension
$48$ the value a bound has to exceed is not a tabulated one but the largest a known
construction gives. That construction
is Edel--Rains--Sloane, evaluated in the dimension itself rather than inherited
from a smaller one (Remark~\ref{rem:ers}). We therefore ran the chain of
\S\ref{sec:ers} in every dimension of Table~\ref{tab:all} above $48$, with the
best published code at each level, and asked of each bound how large the
level-one constant-weight code $A(n,n_1,n_1)$ would have to be for the chain to
overtake it. Divided by the best code available at those parameters, that is the
factor by which the published constructions would have to be beaten.

No bound of Table~\ref{tab:all} falls below the chain, and the factors are far
from uniform. \emph{Dimension $68$ is the
most exposed bound in this paper}: its margin over the chain is $8.9\%$, its
level-one weight is $16$, and $16$ is the one weight at which Edel--Rains--Sloane
built a strong code of their own, giving $A(64,16,16)\ge30\,828$. A code with
$A(68,16,16)\ge45\,235$ would overturn Theorem~\ref{thm:6869} in that dimension,
a factor $1.47$; for calibration, $30\,828$ beats the grid map of
\S\ref{ss:d6263} at the same parameters by $1.88$. No $A(68,16,16)$ has been
published, so the exposure is to a better code and not to the literature as it
stands. The next smallest factor is $4.06$, in dimension $95$; every other
dimension exceeds $5$. Dimensions $62$, $63$ and $96$ are the chain, so the question
does not arise there.

\subsection{The maximum independent set in \texorpdfstring{$\Gamma_\gamma(\Leech)$}{Gamma(Leech)}}\label{ss:class}

Section~\ref{sec:cap} depends on the largest number of minimal lines of $\Leech$
that are pairwise at $|\cos|\le\tfrac14$. The record is $248$ lines
\cite{ma-etal2025}, after $240$ \cite{cohn-jiao-kumar-torquato2011} and $244$
\cite{kallal-kan-wang2017}; the best upper bound is $\lfloor 4680/11\rfloor=425$
\cite[\S7]{cohn-jiao-kumar-torquato2011}. By \eqref{eq:count} one further line in
every independent set contributes $2\sum_i(|Z_i|-1)$, so raising the record by a
single line gives $+2$, $+8$, $+16$, $+32$, $+52$, $+96$ and $+168$ in dimensions
$25,\dots,31$ respectively.

The Delsarte linear programme, the Lov\'asz theta function and the Hoffman ratio
bound on the association scheme whose relations on the $98\,280$ lines are
$|\ip{u}{u'}|\in\{0,1,2\}$ all give $4680/11$, so no two-point relaxation
improves it. The three-point semidefinite bound of Bachoc and
Vallentin \cite{bachoc-vallentin2008} does not improve it at any truncation
computed, the positivity constraints never being active. Closing the gap requires a four-point relaxation
or an argument using the lattice structure.

\subsection{Dimension 17}\label{ss:d17}

The layered family over a $\Lambda_{16}$ equator reduces in dimension $17$ to a
single integer: its largest member has $5\,346+2\alpha$ points, where $5\,346$ is
the kissing number of $\Lambda_{17}$ and $\alpha$ is again an independence number
in a Cayley graph. A Parseval argument excludes the support sizes $7$ and $8$,
leaving the supports of \cite{cohn-li2024}; the remainder of the configuration is
exactly $512$ vectors up to sign.

The value $\alpha=192$ is attained by \cite{cohn-li2024}, and it is optimal. The
vertex set is not an arbitrary $1\,024$-element set: it is
$K'=\{b\in\mathrm{RM}(2,4):|b\cap S_0|\text{ even}\}$, a binary linear code, hence
an elementary abelian group $\mathbb{F}_2^{10}$, and $W_4\subseteq K'$ is symmetric.
The graph is therefore a translation scheme, Delsarte's linear programming bound
applies, and for a Cayley graph on an abelian group that bound is exactly
$\vartheta'$. Its value is $192$.

\begin{proposition}\label{prop:d17}
$\alpha=192$, so the layered family of dimension $17$ contains no configuration
with more than $5\,730$ points.
\end{proposition}

\begin{proof}
Let $C\subseteq K'$ have no two elements differing by a member of $W_4$. For any
$g\colon K'\to\Q$ with $\widehat{g}\ge0$,
\[
  0\;\le\;\tfrac{1}{1024}\sum_{y}\widehat{g}(y)\Bigl|\sum_{c\in C}\chi_y(c)\Bigr|^{2}
   \;=\;\sum_{c,c'\in C}g(c-c')\;\le\;|C|\,g(0)-|C|\bigl(|C|-1\bigr)
\]
whenever $g\le-1$ off $W_4\cup\{0\}$, so $|C|\le 1+g(0)$. Take $g$ whose
normalised Fourier coefficients $b_y=\widehat{g}(y)/1024$ depend only on the
eigenvalue $\lambda(y)=\sum_{w\in W_4}\chi_y(w)$ and on whether $y$ is trivial on
the fibre subgroup. With $b_y=0$ on the remaining $118$ characters:
\setlength{\arraycolsep}{4pt}
\[
\begin{array}{c|cccccc}
(\lambda,\ y|_H\text{ trivial}) & (30,\text{no}) & (12,\text{yes}) & (12,\text{no})
  & (6,\text{no}) & (-4,\text{no}) & (-2,\text{no})\\ \hline
\#\{y\} & 16 & 10 & 120 & 160 & 360 & 240\\
b_y & 65/80 & 40/80 & 30/80 & 28/80 & 10/80 & 9/80\\
\text{mass} & 13 & 5 & 45 & 56 & 45 & 27
\end{array}
\]
\setlength{\arraycolsep}{5pt}
Then $b\ge0$ and $g(0)=\sum_y b_y=191$, while $g\le-1$ at all
$963$ points outside $W_4\cup\{0\}$: exactly $-1$ at $887$ of them, $-\tfrac{21}{5}$
at $60$ and $-9$ at $16$. Scaled by $80$ every quantity is an integer, so the
verification is one Walsh--Hadamard transform in exact arithmetic. This gives
$\alpha\le192$; the configuration of \cite{cohn-li2024} is an independent set of
size $192$ --- six pairwise non-adjacent fibres --- so $\alpha=192$.
\end{proof}

Two coarser relaxations both stop at $256$, and for the same reason. The spectrum
of the graph is $60^{1}30^{16}12^{130}6^{160}(-2)^{240}(-4)^{375}(-10)^{96}
(-20)^{6}$, so the ratio bound is $1024\cdot20/80=256$; and the quotient on the
$32$ fibres has clique number $4$, forcing a cover by eight $4$-cliques of exact
local bound $32$. Both are satisfied by the uniform fractional point $1/4$ on the
fibres, as is every clique, block-pair and odd-hole cut, so no relaxation phrased
on the quotient can beat $256$. The linear programme over the group does not
factor through the quotient. As above, the conclusion applies to the family
alone.

\subsection{Verification}\label{ss:verif}

The package \cite{repo} produces and checks every number above. It applies the
following standards.

\begin{itemize}
\item \emph{No floating-point arithmetic decides any bound of
      Table~\ref{tab:all}.} Every linear programming bound carries an exact
      rational dual certificate (Lemma~\ref{lem:dual}); the solver is used only
      to propose a dual vector. What decides the remaining bounds is an integer
      comparison. Two scans are too large to run in arbitrary-precision
      integers --- the axis layer of dimension $95$ compares $93\,150^2$ inner
      products --- and are run in floating point on integer operands bounded so
      that every product and every partial sum is exactly representable. The
      bound is asserted before the scan; where it does not hold, the operand is
      split into two limbs at bit $31$ and recombined exactly.
\item \emph{Every Gram matrix is proved realizable before use}, by
      Lemma~\ref{lem:realize} or by exhibiting a $k$-tuple that carries it. The
      ones neither licensed nor exhibited are recorded and discarded.
\item \emph{Configurations small enough to write down are checked in
      coordinates.} The dimension-$27$ configuration is verified by two
      deliberately different arguments, one of which compares all
      $\binom{200540}{2}=20\,108\,045\,530$ pairs.
\item \emph{The cross-section method is validated against every independently
      known value} (Table~\ref{tab:validate}). The inequalities and the counting
      formula of the layered construction are re-derived from first principles
      in exact rational arithmetic, independently of the code that applies them.
\item \emph{Large intermediate data is regenerated by deterministic scripts}
      that re-verify every property the claims depend on.
\end{itemize}

\section*{Use of AI assistance}

The construction of \S\ref{ss:d27} was found by Claude Opus~5 (Anthropic, maximum
reasoning effort) in a chat session in which the author provided the problem,
Cohn's table \cite{cohn-table} and Cohn's data set \cite{cohn-data}. It appeared
at the fortieth prompt; most of the thirty-nine earlier prompts were some form of
``try harder''.

The remaining constructions, the package \cite{repo} and this exposition were then
developed by the author with Claude Code (Anthropic), a coding agent built on the
same model. The agent proposed constructions, implemented and ran the verification
code, and drafted the exposition. The author set the direction, selected the lines
to pursue and reviewed the text and the certificates.

In neither setting did the model predict success, and it said so in plain terms:
``I have not produced a new SOTA \dots\ continuing would be me generating
variations faster than either of us can check them.''

\bibliographystyle{amsplain}
\bibliography{kissing46}


\end{document}
